\documentclass[11pt,a4paper]{article}

\usepackage{iftex}
\ifPDFTeX
  \usepackage[T1]{fontenc}
  \usepackage[utf8]{inputenc}
  \usepackage{lmodern}
\else
  \usepackage{fontspec}   % XeTeX (tectonic): Unicode text such as degree signs and accents
\fi
\usepackage[margin=2.3cm]{geometry}
\usepackage{amsmath,amssymb}
\usepackage{booktabs,tabularx,enumitem,xcolor}
\usepackage[colorlinks=true,linkcolor=blue!50!black,urlcolor=blue!50!black]{hyperref}
\newcolumntype{Y}{>{\raggedright\arraybackslash}X}
\newcolumntype{R}[1]{>{\raggedright\arraybackslash}p{#1}}
\Urlmuskip=0mu plus 1mu   % code spans may break between any two characters
\newcommand{\code}[1]{\texttt{\nolinkurl{#1}}}
\newcommand{\codeline}[1]{\texttt{\detokenize{#1}}}   % code with spaces (not for \\ % #)
\newcommand{\finding}[1]{\textcolor{red!60!black}{\textbf{[finding]}}~#1}

% unit shorthands (math mode)
\newcommand{\tha}{\mathrm{t\,ha^{-1}}}
\newcommand{\kgNha}{\mathrm{kg\,N\,ha^{-1}}}
\newcommand{\mxnha}{\mathrm{MXN\,ha^{-1}}}
\newcommand{\tCha}{\mathrm{t\,C\,ha^{-1}}}
\newcommand{\mxnt}{\mathrm{MXN\,t^{-1}}}
\newcommand{\mxnkgN}{\mathrm{MXN\,(kg\,N)^{-1}}}
\newcommand{\CVaR}{\operatorname{CVaR}}
\setlist{itemsep=1pt,topsep=2pt,parsep=0pt}
% tables stay where they are written (no floating) and carry a numbered caption below
\makeatletter
\renewenvironment{table}[1][]{\par\vspace{3pt}\begin{center}\footnotesize}{\end{center}\par\vspace{1pt}}
\renewcommand{\caption}[1]{\par\nopagebreak[4]\vspace{3pt}\refstepcounter{table}\noindent\parbox{0.94\linewidth}{\footnotesize\raggedright\textbf{Table~\thetable.} #1}\par\vspace{2pt}}
\renewcommand\section{\@startsection{section}{1}{\z@}{-2.2ex plus -.5ex minus -.2ex}{1.0ex plus .2ex}{\normalfont\Large\bfseries}}
\renewcommand\subsection{\@startsection{subsection}{2}{\z@}{-1.6ex plus -.4ex minus -.2ex}{0.6ex plus .1ex}{\normalfont\large\bfseries}}
\makeatother
\setlength{\parindent}{0pt}
\setlength{\emergencystretch}{3em}
\setlength{\parskip}{3pt}
\setlength{\abovedisplayskip}{4pt plus 2pt minus 2pt}
\setlength{\belowdisplayskip}{4pt plus 2pt minus 2pt}
\renewcommand{\arraystretch}{1.04}

\title{\vspace{-1.6cm}The NSGA planner: what it computes, line by line\\[0.2em]
\large A typeset walkthrough of the mathematics, the code and the numbers}
\author{Prepared for Isaac by the \texttt{isaacs-inspector} agent \\ \small October 26, 2026}
\date{}

\begin{document}
\maketitle
\vspace{-1.1cm}

\section{Takeaways and what changed after review}
\label{sec:take}

\noindent\fcolorbox{blue!50!black}{blue!4}{\begin{minipage}{\dimexpr\linewidth-2\fboxsep-2\fboxrule\relax}
\small\textbf{If you read only this box.}
\begin{enumerate}[leftmargin=1.6em,itemsep=2pt,topsep=2pt]
\item \textbf{What it is.} Scores $4{,}284$ strategies $x=(\omega,s,r)$ (rotation, fertilizer scale, stover retention) on five objectives by interpolating $2{,}016$ WOFOST
  crop-simulation seasons; NSGA-II/III (evolutionary multi-objective search) is checked against brute force over all $4{,}284$.
\item \textbf{The arithmetic is right.} An independent pandas recomputation of 10 strategies matches the planner to a relative $1.2\times10^{-16}$; all 35 unit tests pass.
\item \textbf{Two inputs set the conclusions, not the algorithms:} how much yield nitrogen buys (calibrated WOFOST: none, $2.301\to2.312\,\tha$; default: almost linear) and placeholder prices (the faba bean price used is $3.3\times$ SIAP's latest value).
\item \textbf{``N lost'' restates the dose:} $\approx0.49\,D$ in the default world (by the fixed 50\% recovery constant), $\approx0.95\,D$ in the calibrated one. It is not a measured loss.
\item \textbf{The big fronts are partly artifacts:} $483$ calibrated strategies come from dominance resistance (Section~\ref{sec:dominance}). A float-noise bug had kept $21$ more on the front ($504$ before) and produced one of the picks; it is fixed (Section~\ref{sec:artifact}).
\item \textbf{NSGA-III never ``plateaued'' at 89--94\%:} that was a reporting bug (\code{result.X}); the final population covers $95$--$98.5\%$ of the hypervolume, and NSGA saves no evaluations over brute force here.
\item \textbf{Chebyshev picks and the robust ``hedge dose'' are unstable.} The hedge ($48$--$132\,\kgNha$) holds only under range-normalized regret ($176$--$200$ in absolute pesos).
\item \textbf{Resilience is weak evidence.} Bad-year income is the mean of the worst 4 of 16 years, which are frost-truncated seasons (correlation with yield $-0.70$), not droughts; the holdout cannot test
  drought; the SMAP agreement ($r=0.87$) feeds nothing and could not be re-run.
\end{enumerate}
\end{minipage}}

\noindent\textbf{Scope.} This document explains \code{src/crop_models/nsga_planner.py} (multi-objective planning of rainfed rotations at Texcoco, Mexico) as it stands in the working tree on
October 26, 2026 (HEAD \code{21e30ec} plus uncommitted changes), replaces the plain-text draft \code{src/crop_models/NSGA_WALKTHROUGH.md} (keeping its structure and status table) and is meant to be read alone.
Every number was recomputed from \code{data/nsga/grid.csv} or the saved fronts in \code{data/nsga/*/} with the repository's own code (Section~\ref{sec:verify}).
Statements are labeled \emph{code} (what the code does), \emph{docs} (what repository documents say) or \emph{my reading} (interpretation).

\newpage
\subsection*{Status of the earlier review findings}

The first review read the code just before these changes. ``Fixed'' means I confirmed the fix in the current code or data.

\begin{table}[h]
\begin{tabularx}{\linewidth}{@{}R{0.40\linewidth}Y@{}}
\toprule
Finding & Status \\
\midrule
NSGA-III's ``89--94\% plateau'' is a reporting artifact: pymoo's \code{result.X} keeps one strategy per occupied reference direction.
& \textbf{Fixed.} \code{optimize} returns the non-dominated set of the final population \code{result.pop} (\code{nsga_planner.py:323--326}; test \code{tests/test_nsga_planner.py:199}); coverage $95$--$98.5\%$ (Table~\ref{tab:bug}). The ``thin front'' explanation is removed. \\
``Worst 20\% of years'' is the worst 25\% for 16 years ($q=\lceil0.2\cdot16\rceil=4$).
& \textbf{Fixed by relabeling} to the worst quarter, \codeline{CVAR_SHARE = 0.25} (\code{nsga_planner.py:82}). No number changes. \\
Default-world stover and roots are not capped like grain.
& \textbf{Documented}, not changed (shared simulator code); still open (Section~\ref{sec:objectives}). \\
Doc slips (``3\,kg'' should be 4; Reproduce commands did not match \code{run.json}).
& \textbf{Fixed.} The saved runs reproduce here (Section~\ref{sec:nsga}). \\
Latest SIAP faba bean price (2025: $5{,}030\,\mxnt$) is $0.30\times$ the price used.
& \textbf{Documented.} The price used is the median of $16{,}548$, $16{,}524$, $5{,}030$. \\
The minimax dose depends on how regret is normalized ($176$--$200\,\kgNha$ in absolute pesos).
& \textbf{Verified and documented}; derived in closed form in Section~\ref{sec:robust}. \\
``Benign years'' for 2022--2025 is unsupported; ``N loss is a lower bound'' unverifiable; IGBP's carbon stock is read by no objective.
& \textbf{Rewritten} in the results document. \\
Calibrated picks sit inside the first interpolation cell $[0,0.25]$, whose interior was never simulated.
& \textbf{Documented}; a finer low-dose grid is a next step. After the float fix the calibrated picks are at $s=0$ (a simulated node), but the $231$ maize front members still lie inside the cell. \\
$21$ calibrated front members existed only through a $4.4\times10^{-16}\,\tha$ float difference (second review); the first review's ``no pick changes'' was wrong.
& \textbf{Fixed.} \code{pareto_front} (\code{nsga_planner.py:231--244}) compares objectives rounded to $9$ decimals (\codeline{FRONT_DECIMALS = 9}; tests \code{tests/test_nsga_planner.py:163,172}). Front sizes, the coverage tables and the calibrated sustainability-first pick changed (Section~\ref{sec:artifact}). \\
\bottomrule
\end{tabularx}
\caption{Findings of the first review and their status.}
\label{tab:status}
\end{table}

\subsection*{Corrections made while re-typesetting this document}

Every number of the Markdown draft was re-verified; six changed (marked where they occur).

\begin{table}[h]
\begin{tabularx}{\linewidth}{@{}R{0.26\linewidth}R{0.24\linewidth}Y@{}}
\toprule
Item & Markdown draft & Re-verified value \\
\midrule
Maize--faba bean 2018 N-loss example & ``the faba phase took up 97.5 kg'' & It is the \emph{maize} phase ($Q=97.52$, $Q(0)=88.54$, $D=80$); the faba phase has $D=20$ and no extra uptake. \\
Front members dominated once yield is dropped & $374$ of $462$ (leaving $88$) & $444$ of $462$ (leaving $18$) \\
Standing mineral N versus ``lost'' N (calibrated, maize, $s=1$) & ``about 3\% of the 152 kg lost'' & $68.7\,\kgNha$ more mineral N at the start of the 2021 season than at $s=0$, i.e.\ $2.8\%$ of the 16-year cumulative ``loss'' $16\times152.2=2{,}435\,\kgNha$ \\
Crossing of the two income-shortfall lines & $N^\ast=0.180/(a+b)=87.5$ & $87.3$ by the closed form; $89.4$ exact on the piecewise-linear curves; lattice optimum $88$ \\
Held-out net income of the two picks & $11{,}598$ vs $6{,}086$ (picks of an earlier front) & $11{,}598$ vs $6{,}158$ for the picks on the current saved front \\
Faba bean area behind the price & ``4--12 ha'' & $10$--$12$ ha of faba bean (the 4--12 range is barley and faba bean together) \\
\bottomrule
\end{tabularx}
\caption{Numbers corrected relative to \code{NSGA_WALKTHROUGH.md}.}
\label{tab:corr}
\end{table}

\section{The picture and the pipeline}
\label{sec:pipeline}

\textbf{The question.} A farmer at Texcoco (Valley of Mexico, $19.53^\circ$N, $98.85^\circ$W, rainfed, one crop per
rainy season) can choose a rotation, a fertilizer level and how much stover (the leaves and stems left after harvest)
to leave in the field. Which choices are good \emph{jointly} for yield, income, nitrogen lost to the environment, soil
carbon returned and surviving bad years? There is no single best answer, so the planner returns the
\emph{Pareto front}: the strategies that no other strategy beats on every objective at once.

\textbf{What goes in and out.} In: a table of $2{,}016$ simulated crop seasons (\code{data/nsga/grid.csv}) from WOFOST
(WOrld FOod STudies, a daily crop growth model; run through PCSE, the Python Crop Simulation Environment) for two parameter sets, plus prices
and costs. Out: for every strategy five numbers, the Pareto front of those numbers, and (optionally) one compromise pick.
Where it sits: it is the last stage of the pipeline, downstream of ingest, simulation and calibration.

\begin{center}\small
\begin{tabular}{@{}c@{\,}c@{\,}c@{\,}c@{\,}c@{\,}c@{\,}c@{}}
NASA data & $\rightarrow$ & WOFOST runs & $\rightarrow$ & \code{grid.csv} & $\rightarrow$ & \code{Grid} (arrays) \\[2pt]
 & & & & & & $\downarrow$ \\[2pt]
\code{choose} / \code{RobustPlanner} & $\leftarrow$ & NSGA-II/III + brute force & $\leftarrow$ & \code{Planner.objectives} & $\leftarrow$ & interpolation along $s$
\end{tabular}
\end{center}

The planner chain in code: \code{Grid} (\code{nsga_planner.py:106}) $\to$ \code{Planner.annual} (:164) $\to$
\code{Planner.objectives} (:182) $\to$ \code{brute_force} (:258), \code{optimize} (:266), \code{front_quality} (:329),
\code{choose} (:348), \code{RobustPlanner} (:190).

\subsection{Which dataset feeds which stage}

\begin{table}[h]
\begin{tabularx}{\linewidth}{@{}R{0.22\linewidth}Y R{0.24\linewidth}@{}}
\toprule
Dataset & What it provides here & Where used \\
\midrule
POWER (NASA Prediction Of Worldwide Energy Resources) & Daily weather 2005--2025 at the site: radiation, temperature, rain. Drives every simulated season; its year-to-year spread is the only source of risk in the bad-year objective. & \code{wofost_rotations.py:342} \\
IGBP-DIS (soil data of the International Geosphere--Biosphere Programme, via NASA Earthdata) & Field capacity $0.373$ and wilting point $0.137\ \mathrm{m^3\,m^{-3}}$ (0--100\,cm); total soil N $12{,}709\,\kgNha$, of which $0.5\%$ per season is mineralized: $0.005\times12{,}708.6=63.5\,\kgNha$ & \code{wofost_rotations.py:149-150,157,118} \\
GLDAS (Global Land Data Assimilation System) & Soil porosity $0.437\ \mathrm{m^3\,m^{-3}}$ & \code{wofost_rotations.py:151} \\
IGBP-DIS soil carbon & $130.4\ \mathrm{t\,C\,ha^{-1}}$ stock, written to \code{grid.json} only. \emph{No objective reads it.} & \code{nsga_planner.py:526-530} \\
SMAP (Soil Moisture Active Passive) & An earlier validation of simulated soil water (below); \emph{not used by the planner}. & \code{wofost_rotations.py:391} \\
\midrule
SIAP (Mexican agriculture ministry; not NASA) & Prices (median of the last three priced years: maize $7{,}981$, barley $5{,}240$, faba bean $16{,}524\,\mxnt$) and the observed yields behind the calibration. & \code{farm_tool.py:98} \\
Placeholders (no source) & Other costs ($14{,}000$, $9{,}000$, $11{,}000\,\mxnha$), N price $25\,\mxnkgN$, stover price $1{,}000\,\mxnt$, carbon fraction $0.45$. & \code{nsga_planner.py:78-81} \\
\bottomrule
\end{tabularx}
\caption{Inputs of the planner and their provenance. The grid has 8 rotation-phases ($1+2+2+3$) $\times$ 21 years
$\times$ 6 fertilizer scales $\times$ 2 parameter sets $=2{,}016$ rows; the planner uses 2006--2021 (design) and 2022--2025 (holdout).}
\end{table}

\subsection{The SMAP check: what it is and what it is not}

\emph{Code} (\code{wofost_rotations.py:391--402}): the simulated 0--100\,cm soil water $\mathrm{WWLOW}$ (cm) is divided by the rooting depth
of $100$\,cm to give a volumetric fraction $\mathrm{SM}^{\mathrm{sim}}_d$ ($\mathrm{m^3\,m^{-3}}$), and compared with SMAP's L4 root-zone moisture
$\mathrm{SM}^{\mathrm{SMAP}}_d$ (daily means of 3-hourly values) by a Pearson correlation over the overlapping days $d$,
\begin{equation}
r=\operatorname{corr}_d\!\big(\mathrm{SM}^{\mathrm{sim}}_d,\ \mathrm{SM}^{\mathrm{SMAP}}_d\big),\qquad
\mathrm{SM}^{\mathrm{sim}}_d=\frac{\mathrm{WWLOW}_d}{100\ \mathrm{cm}}.
\label{eq:smap}
\end{equation}
\emph{Docs}: $r=0.87$ over $731$ days (2023--2024) for the maize-every-year rotation
(\code{docs/technical_report/tables/smap.tex}). \textbf{I could not re-run this}: it needs a PCSE simulation, which was
excluded here; I also cannot tell from the repository which parameter set (default or calibrated) produced it.

What it is not (\emph{my reading}): it calibrates nothing, and \code{nsga_planner.py} never reads SMAP. Pearson $r$ ignores bias and scale. The $731$ daily values are not
independent (the lag-1 autocorrelation of the SMAP series in \code{data/smap_texcoco_2023_2024.csv} is $0.998$, lag 30 is $0.83$; I computed this), so $r$ is dominated by the shared
seasonal wet--dry cycle and the effective sample is far smaller than $731$. And it compares one simulated field with a $\sim9$\,km SMAP cell. It says the simulated water follows the
right seasons, not that drought years are reproduced.

\subsection{Notation used throughout}

\begin{table}[h]
\begin{tabularx}{\linewidth}{@{}R{0.15\linewidth}Y R{0.27\linewidth}@{}}
\toprule
Symbol & Meaning & Unit or value \\
\midrule
$\omega$ & rotation (categorical): maize, maize--barley, maize--faba bean, maize--barley--faba bean & --- \\
$s$ & fertilizer scale, a multiple of each crop's dose ($s=1$: maize 160, barley 60, faba bean 40) & dimensionless, $[0,1.25]$ \\
$r$ & stover retention, the share of stover left in the field & fraction, $[0,1]$ \\
$x=(\omega,s,r)$ & a strategy; the lattice is $\Omega$, $|\Omega|=4{,}284$ & --- \\
$N=160\,s$ & ``maize N'', the maize dose implied by $s$ & $\kgNha$ \\
$t,\ T$; $u,\ n_u$ & design year ($2006$--$2021$, $T=16$); phase of the rotation ($n_u=1,2,2,3$) & --- \\
$c_{tu}$ & crop grown in year $t$ in phase $u$ & maize, barley, faba bean \\
$G,\,D,\,Q,\,S,\,B$ & grain, N dose, N uptake, stover (leaves incl.\ dead + stems), roots, per season & $\tha$ dry matter, $\kgNha$, $\kgNha$, $\tha$, $\tha$ \\
$\pi_c,\,K_c$; $\pi_N,\,\pi_S,\,\phi$ & grain price, other cost of crop $c$; N price $25$, stover price $1{,}000$, carbon fraction $0.45$ & $\mxnt$, $\mxnha$; $\mxnkgN$, $\mxnt$, $\mathrm{t\,C\,(t\,DM)^{-1}}$ \\
$Y_t,V_t,L_t,C_t$ & annual (phase-averaged) yield, net income, N loss, carbon input & $\tha$, $\mxnha$, $\kgNha$, $\tCha$ \\
$f_1,\dots,f_M$; $\sigma_i$ & objectives in minimization form $f_i=\sigma_i\times$ natural value, $\sigma_i=\pm1$; $M\in\{2,\dots,5\}$ & --- \\
$\alpha,\ q$ & tail share $0.25$ and tail count $q=\lceil\alpha T\rceil=4$ & --- \\
$\theta$ & WOFOST parameter set (``world''): default (def) or calibrated (cal) & --- \\
\bottomrule
\end{tabularx}
\caption{Symbols. A symbol keeps its meaning for the whole document.}
\end{table}

\section{The grid and linear interpolation along the fertilizer scale}
\label{sec:grid}

\textbf{Plain language.} WOFOST was run only at six fertilizer levels. A strategy at a finer level, say $s=0.6$, is
scored by blending the two neighboring simulated levels, season by season, as if yield were a straight line between
them. This is what makes a strategy cost about $60\,\mu$s instead of a 3--10\,s simulation.

\subsection{The equation}

The grid nodes are $s_0,\dots,s_5=0,\,0.25,\,0.5,\,0.75,\,1,\,1.25$ (\code{GRID_SCALES}, \code{nsga_planner.py:60}).
For any simulated quantity $X\in\{G,D,Q,S,B\}$ of one (year, phase) cell, a strategy with scale $s$ (clipped to
$[0,1.25]$) uses
\begin{equation}
j(s)=\min\!\big\{4,\ \max\{j:\ s_j\le s\}\big\},\qquad
w=\frac{s-s_{j}}{s_{j+1}-s_{j}},\qquad
X(s)=(1-w)\,X_{j}+w\,X_{j+1}.
\label{eq:interp}
\end{equation}
Here $w\in[0,1]$ is the weight on the upper node (dimensionless) and $X_j$ is the value simulated at node $s_j$.
Code: \code{Grid.interpolate}, \code{nsga_planner.py:136--142}; the clip is line 138, $j$ is line 139, $w$ is line 140
and the blend is line 141. At a node $w=0$ and the table value is returned exactly (test \code{tests/test_nsga_planner.py:63}).

\textbf{Worked number} (default parameters, pure maize, 2006, $s=0.6$). Here $j=2$ ($s_2=0.5$) and
\begin{equation}
\begin{aligned}
w &= \frac{0.6-0.5}{0.75-0.5}=0.4,\\
G(0.6) &= 0.6\times5.965005+0.4\times7.993159 = 3.579003+3.197264 = 6.776267\ \tha .
\end{aligned}
\end{equation}
The code returns $6.776266824$; hand and code agree.

\subsection{What the nodes look like}

\begin{table}[h]
\begin{tabular}{@{}llcccccc@{}}
\toprule
 & & \multicolumn{6}{c}{fertilizer scale $s_j$ (maize dose $D=160\,s_j$)} \\
\cmidrule(l){3-8}
Parameters & Mean over 2006--2021, pure maize & $0$ & $0.25$ & $0.5$ & $0.75$ & $1$ & $1.25$ \\
\midrule
default & yield $\bar G$ [$\mathrm{t\,ha^{-1}}$] & $1.918$ & $3.038$ & $4.162$ & $5.252$ & $6.482$ & $7.599$ \\
default & uptake $\bar Q$ [$\mathrm{kg\,N\,ha^{-1}}$] & $66.6$ & $86.8$ & $107.0$ & $127.2$ & $147.4$ & $167.6$ \\
calibrated & yield $\bar G$ [$\mathrm{t\,ha^{-1}}$] & $2.301$ & $2.312$ & $2.312$ & $2.312$ & $2.312$ & $2.312$ \\
calibrated & uptake $\bar Q$ [$\mathrm{kg\,N\,ha^{-1}}$] & $59.9$ & $67.3$ & $67.8$ & $67.8$ & $67.8$ & $67.8$ \\
\bottomrule
\end{tabular}
\caption{Node means of grain yield $\bar G$ and N uptake $\bar Q$ (pure-maize rotation, design years).}
\label{tab:nodes}
\end{table}

\textbf{Default world: a near-linear nitrogen response.} Uptake rises $20.2\,\kgNha$ per $40\,\kgNha$ of dose, an
\emph{apparent recovery}
\begin{equation}
\mathrm{AR}(s)=\frac{\bar Q(s)-\bar Q(0)}{\bar D(s)}\quad[\text{dimensionless}],\qquad
\mathrm{AR}_{\mathrm{def}}(0.25)=\frac{86.816-66.584}{40}=0.506,
\label{eq:ar}
\end{equation}
which is the hand-set \codeline{N_RECOVERY = 0.5} (\code{wofost_rotations.py:117}, ``share of fertilizer N the crop can access'')
plus a little carry-over. Yield rises by about $1.1\,\tha$ per node.

\textbf{Calibrated world: no response, and why.} Yield moves from $2.3008$ to $2.3125\,\tha$ ($+0.0117$, $+0.51\%$) and
uptake saturates at $67.75\,\kgNha$ from $s=0.5$ on, giving $\mathrm{AR}_{\mathrm{cal}}(1)=(67.75-59.93)/160=0.049$.
\emph{Docs} (\code{CALIBRATION_RESULTS.md}): the calibration absorbed the Texcoco yield gap into photosynthesis (maize scaled
by $0.47$, grid.json), so growth is limited by carbon rather than nitrogen. The crop's N demand tops out near
$68\,\kgNha$ while the soil already supplies about $63.5$ (mineralization above) plus carry-over, so extra fertilizer has
nothing to do. \emph{My reading}: this is a property of the fitted parameter set, whose only nitrogen evidence is one
station trial (unfertilized/fertilized ratio $0.93$; the model's is $0.995$), not an agronomic finding.
\emph{Docs} on the fit: the calibrated knobs per crop are a thermal-time (phenology) scale and a photosynthesis scale (maize $0.8$ and $0.47$, \code{grid.json}; the $0.8$ sits on the agronomic lower bound), with N recovery and soil mineralization left at their defaults; they were fit to the 16 yearly SIAP yields of 2006--2021, and 2022--2025 was held out.

\subsection{Assumptions and where they break}

\begin{itemize}
\item \textbf{Histories are mixed.} Each node is a 21-year simulation at a fixed $s$, chaining water and soil N from season
  to season. A season at $s=0.6$ would carry a $0.6$-history; Equation~\eqref{eq:interp} blends the $0.5$- and
  $0.75$-histories instead.
\item \textbf{Straight line between nodes.} Reasonable where the response is near-linear (default), not where it jumps
  (calibrated uptake between $s=0$ and $0.25$, a maximum change of $29\,\kgNha$ in one season).
\item \textbf{The lattice is $10\times$ finer than the nodes.} Strategies sit every $0.025$, nodes every $0.25$. The calibrated
  picks are at $s=0$ (a simulated node), but the $231$ maize strategies on the five-objective front ($s=0$--$0.25$) rest on interpolation inside the first cell $[0,0.25]$, whose interior was never simulated.
\end{itemize}

\textbf{Leave-one-node-out error.} Predict each interior node $j=1,\dots,4$ from its two neighbors,
$\hat X_j=\tfrac12(X_{j-1}+X_{j+1})$, and compare with the simulated $X_j$ over all rotations, years and phases. The
spacing is twice the real one, so this overstates the error of Equation~\eqref{eq:interp}.

\begin{table}[h]
\begin{tabular}{@{}llccc@{}}
\toprule
World & Quantity & mean $|\hat X_j-X_j|$ per season & 95th pct & max \\
\midrule
default & grain yield ($\tha$) & $0.048$ & $0.20$ & $0.82$ \\
default & N uptake ($\kgNha$) & $0.014$ & $0.043$ & $0.42$ \\
calibrated & grain yield ($\tha$) & $0.000$ & $0.000$ & $0.064$ \\
calibrated & N uptake ($\kgNha$) & $0.24$ & $0.000$ & $13.2$ \\
\bottomrule
\end{tabular}
\caption{Leave-one-node-out error. The error of the 16-year, all-phase \emph{mean} is at most $0.069\,\tha$ for default
yield (about $1\%$) and $3.42\,\kgNha$ for calibrated uptake.}
\label{tab:loo}
\end{table}

\section{The five objectives, and one strategy worked end to end}
\label{sec:objectives}

\textbf{Plain language.} For a given strategy the planner builds, for each of the 16 design years, four farm-level numbers
(yield, net income, N lost, carbon input), averaging over the rotation's phases (a hectare is, on average, half maize
and half faba bean each year in a two-crop rotation). The objectives are the 16-year means of those, plus a risk
measure: the mean net income of the worst quarter of years.

\subsection{The equations}

For one (year, phase) cell, let $G,D,Q,S,B$ be the interpolated values from Equation~\eqref{eq:interp} at the strategy's
$s$, and $Q(0)$ the N uptake of the \emph{same cell} at the zero-fertilizer node. With crop $c=c_{tu}$ in that cell, the
four annual numbers are phase averages ($n_u$ phases):
\begin{align}
Y_t &= \frac1{n_u}\sum_{u=1}^{n_u} G_{tu}
  && [\tha\ \text{dry matter grain}] \label{eq:yield}\\
V_t &= \frac1{n_u}\sum_{u=1}^{n_u}\Big[\pi_c\,G_{tu}-\big(K_c+\pi_N D_{tu}\big)+\pi_S\,(1-r)\,S_{tu}\Big]
  && [\mxnha] \label{eq:net}\\
L_t &= \frac1{n_u}\sum_{u=1}^{n_u}\max\!\Big\{0,\ D_{tu}-\big(Q_{tu}-Q_{tu}(0)\big)\Big\}
  && [\kgNha] \label{eq:nloss}\\
C_t &= \frac1{n_u}\sum_{u=1}^{n_u}\phi\,\big(B_{tu}+r\,S_{tu}\big)
  && [\tCha] \label{eq:carbon}
\end{align}
The five objectives (all in natural units; the search minimizes $f_i=\sigma_i\times$ objective with $\sigma=-1$ for
the maximized ones, \code{NATURAL_SIGN}, :69):
\begin{align}
\bar Y&=\frac1T\sum_{t=1}^T Y_t, \qquad \bar V=\frac1T\sum_{t=1}^T V_t, \qquad
\bar L=\frac1T\sum_{t=1}^T L_t, \qquad \bar C=\frac1T\sum_{t=1}^T C_t,
\label{eq:means}
\end{align}
\begin{equation}
\CVaR_\alpha=\frac1q\sum_{l=1}^{q}V_{(l)},\qquad
q=\max\{1,\lceil\alpha T\rceil\}=\lceil0.25\times16\rceil=4,\qquad
V_{(1)}\le V_{(2)}\le\cdots\le V_{(T)} .
\label{eq:cvar}
\end{equation}
$\CVaR_\alpha$ (conditional value at risk, here of a gain, i.e.\ the expected shortfall of income) is the mean of the $q$
lowest annual incomes $V_{(l)}$ (order statistics, $\mxnha$). Because $\alpha T=4$ is an integer, this is exactly the
textbook empirical estimator of the lower-tail mean at level $\alpha$; for non-integer $\alpha T$ the ceiling makes the tail
slightly larger than $\alpha$.

\begin{table}[h]
\begin{tabularx}{\linewidth}{@{}l Y l@{}}
\toprule
Term & Code & Line \\
\midrule
$G,D,Q,S,B$ at scale $s$ & \codeline{grain, dose, uptake, stover, roots = (g.interpolate(rotation, name, n_scale) for name in FIELDS)} & :174 \\
$\max\{0,\,D-(Q-Q(0))\}$ & \codeline{np.maximum(0.0, dose - (uptake - g.rotations[rotation]["n_uptake_kg_ha"][0]))}; node index 0 is $s=0$ & :175 \\
$V$ & \codeline{grain*price - (cost + e.n_price*dose) + e.stover_price*(1-retention)*stover} & :176--177 \\
$C$ & \codeline{e.carbon_fraction*(roots + retention*stover)} & :178 \\
phase average $\frac1{n_u}\sum_u$ & \code{.mean(axis=1)} (axis 1 = phase), after the floor and the prices are applied cell by cell & :179--180 \\
year mean $\frac1T\sum_t$ & \code{.mean()} in \code{objectives} & :185--186 \\
$\CVaR_\alpha$ & \codeline{tail_mean(values, share=CVAR_SHARE)}: \codeline{k = max(1, int(np.ceil(share*len(values))))}; mean of \code{np.sort(values)[:k]} & :145--148, 187 \\
$\pi_N,\pi_S,\phi,K_c$ & \code{PLACEHOLDER_*}, \code{CARBON_FRACTION}, \code{CVAR_SHARE} & :78--82 \\
\bottomrule
\end{tabularx}
\caption{Symbols to code in \code{nsga_planner.py}.}
\end{table}

\textbf{Assumptions.} The crop in each (year, phase) cell is the same at every $s$ (checked at construction, :131--132).
Retention $r$ is pure accounting: it moves stover between income ($\pi_S$) and carbon ($\phi$) and changes nothing in the
simulation. Prices are constants (no price series), and a phase-averaged hectare is a farm of equal-sized fields.

\subsection{One strategy, end to end}

Strategy: \emph{calibrated} parameters, maize--faba bean, $s=0.5$ (doses $D=80$ for maize and $20$ for faba bean, so
$50\,\kgNha$ phase average), $r=0.5$, design years 2006--2021. At $s=0.5$, $w=0$ in Equation~\eqref{eq:interp}, so the
node values are used directly. Rotations are offset in time, so in odd years phase~1 is maize and phase~2 faba bean, and
vice versa.

\textbf{Step 1: one year by hand, 2007} (phase~1 is maize). Prices: maize $7{,}981.46$, faba bean $16{,}524\,\mxnt$.

\begin{table}[h]
\setlength{\tabcolsep}{4pt}
\begin{tabular}{@{}lccc@{}}
\toprule
 & maize ($u=1$) & faba bean ($u=2$) & phase average \\
\midrule
$G$ ($\tha$) & $2.1785$ & $1.7079$ & $1.9432$ \\
$D$ ($\kgNha$) & $80$ & $20$ & $50$ \\
$Q=Q(0)$ ($\kgNha$) & $61.936$ & $21.095$ & \\
$S$, $B$ ($\tha$) & $1.3774$, $0.2170$ & $0.8111$, $0.2299$ & \\
\midrule
revenue $\pi_c G$ ($\mxnha$) & $2.1785\times7981.46=17{,}387.5$ & $1.7079\times16524=28{,}221.7$ & \\
other cost $K_c$ & $-14{,}000$ & $-11{,}000$ & \\
N cost $\pi_N D$ & $-25\times80=-2{,}000$ & $-25\times20=-500$ & \\
stover income $\pi_S(1-r)S$ & $1000\times0.5\times1.3774=688.7$ & $1000\times0.5\times0.8111=405.6$ & \\
net $V$ & $2{,}076.2$ & $17{,}127.3$ & $9{,}601.7$ \\
N loss $\max\{0,80-0\}$, $\max\{0,20-0\}$ & $80$ & $20$ & $50.0$ \\
carbon $0.45\,(B+0.5\,S)$ & $0.4075$ & $0.2860$ & $0.3467$ \\
\bottomrule
\end{tabular}
\caption{The year 2007 by hand. The code returns $V=9{,}601.73$, $L=50.0$, $C=0.34675$, $Y=1.94321$ for this year (identical).}
\label{tab:y2007}
\end{table}

\textbf{Step 2: all 16 years.} Every year has dose $50\,\kgNha$; $L_t=50$ except 2018.

\begin{table}[h]
\begin{tabular}{@{}lcrrrr@{}}
\toprule
Year $t$ & phase-1 crop & $Y_t$ ($\tha$) & $V_t$ ($\mxnha$) & $L_t$ ($\kgNha$) & $C_t$ ($\tCha$) \\
\midrule
2006 & \code{fababean} & $2.134$ & $10{,}590$ & $50.0$ & $0.356$ \\
2007 & \code{maize} & $1.943$ & $9{,}602$ & $50.0$ & $0.347$ \\
2008 & \code{fababean} & $1.692$ & $7{,}547$ & $50.0$ & $0.319$ \\
2009 & \code{maize} & $2.413$ & $13{,}499$ & $50.0$ & $0.401$ \\
2010 & \code{fababean} & $0.887$ & $-1{,}727$ & $50.0$ & $0.256$ \\
2011 & \code{maize} & $1.212$ & $3{,}695$ & $50.0$ & $0.376$ \\
2012 & \code{fababean} & $1.553$ & $6{,}937$ & $50.0$ & $0.323$ \\
2013 & \code{maize} & $2.191$ & $12{,}279$ & $50.0$ & $0.366$ \\
2014 & \code{fababean} & $2.252$ & $13{,}213$ & $50.0$ & $0.384$ \\
2015 & \code{maize} & $2.526$ & $15{,}015$ & $50.0$ & $0.427$ \\
2016 & \code{fababean} & $2.388$ & $13{,}961$ & $50.0$ & $0.428$ \\
2017 & \code{maize} & $2.081$ & $10{,}869$ & $50.0$ & $0.384$ \\
2018 & \code{fababean} & $2.724$ & $17{,}375$ & $45.5$ & $0.480$ \\
2019 & \code{maize} & $2.570$ & $15{,}494$ & $50.0$ & $0.422$ \\
2020 & \code{fababean} & $1.676$ & $8{,}006$ & $50.0$ & $0.364$ \\
2021 & \code{maize} & $2.429$ & $14{,}570$ & $50.0$ & $0.391$ \\
\midrule
mean & & $2.042$ & $10{,}683$ & $49.72$ & $0.3765$ \\
\bottomrule
\end{tabular}
\caption{Annual values for the worked strategy (code output; the 2007 row matches Table~\ref{tab:y2007}).}
\label{tab:years}
\end{table}

\textbf{Step 3: the five numbers.}
\begin{equation}
\begin{aligned}
\bar Y &= 2.042\ \tha, \qquad \bar C = 0.3765\ \tCha,\\
\bar L &= \frac{15\times50+45.512}{16}=\frac{795.512}{16}=49.72\ \kgNha,\\
\bar V &= \frac{170{,}925.2}{16}=10{,}682.8\ \mxnha,\\
\CVaR_{0.25} &= \frac{V_{2010}+V_{2011}+V_{2012}+V_{2008}}{4}\\
 &=\frac{-1{,}726.85+3{,}695.48+6{,}937.05+7{,}546.86}{4}=\frac{16{,}452.54}{4}=4{,}113.14\ \mxnha .
\end{aligned}
\label{eq:worked}
\end{equation}
The code's \codeline{Planner.objectives("maize-fababean", 0.5, 0.5)} returns exactly these: yield $2.041985$, net $10{,}682.826$,
N loss $49.7195$, carbon $0.376539$, bad-year income $4{,}113.135$. The worst four years are 2010 ($-1{,}727$), 2011, 2012 and 2008.

\subsection{The difference method for N loss, and what it over- or understates}

\emph{Code.} Equation~\eqref{eq:nloss} says: loss $=$ dose minus the \emph{extra} uptake caused by the fertilizer, where the
no-fertilizer uptake $Q(0)$ is the same year and phase from the $s=0$ node. In 2018 the \emph{maize} phase took up
$97.52$ against $88.54$ unfertilized, an extra $8.98$, so its loss was $80-8.98=71.02$ and the year's average is
$(71.02+20)/2=45.51$. In every other year the uptake is identical with and without fertilizer, so the loss equals the dose.
(\textbf{Correction}: the Markdown draft attributed this to the faba bean phase.)

\emph{My reading} of what the figure measures:
\begin{itemize}
\item \textbf{Default world: a restatement of an assumption.} The simulator lets the crop reach only $50\%$ of applied N
  (\codeline{N_RECOVERY = 0.5}); the crop then takes up essentially all of that half. So $L\approx(1-\mathrm{AR})\,D\approx0.495\,D$:
  maize at $s=1$ loses $79.17$ of $160\,\kgNha$ ($0.495$).
\item \textbf{Calibrated world: the dose, minus a few kilograms.} Maize at $s=1$ loses $152.18$ of $160$ ($0.951$). Decomposition: the
  inaccessible half ($80$, lost by assumption) plus the accessible half the carbon-limited crop does not use
  ($80-7.82=72.18$): $80+72.18=152.18$ (my arithmetic). Part of the unused half is not lost but stays as mineral N: at the start
  of the 2021 season there are $68.7\,\kgNha$ more than at $s=0$, i.e.\ $2.8\%$ of the 16-year cumulative $2{,}435\,\kgNha$.
\item \textbf{Fixed fractions, not weather.} The $0.5$ accessible fraction and the $0.5$ carry-over loss
  (\code{N_CARRYOVER_LOSS}, :120) are constants. Leaching, volatilization, denitrification and N$_2$O are not
  modelled separately, so the figure ranks strategies by dose and cannot be read as an environmental loss.
\item \textbf{Sign unknown.} The first draft called it a lower bound; its sign relative to reality cannot be established from
  this code, and the claim was removed.
\item \textbf{The floor never binds.} $\max\{0,\cdot\}$ is active in $0$ of $6{,}528$ (phase, year, lattice scale) cells per
  world (and in none of the $1{,}536$ raw grid nodes).
\end{itemize}

\subsection{The retention trade-off}

Differentiating Equations~\eqref{eq:net} and \eqref{eq:carbon} in $r$, with $\bar S=\frac1{Tn_u}\sum_{t,u}S_{tu}$ the mean stover,
\begin{equation}
\frac{\partial\bar V}{\partial r}=-\pi_S\,\bar S,\qquad
\frac{\partial\bar C}{\partial r}=+\phi\,\bar S,\qquad
\frac{\partial\bar V/\partial r}{\partial\bar C/\partial r}=-\frac{\pi_S}{\phi}=-\frac{1000}{0.45}=-2{,}222\ \mathrm{MXN\,(t\,C)^{-1}}.
\label{eq:carbonprice}
\end{equation}
Both objectives are exactly linear in $r$ and the implicit carbon price $\pi_S/\phi$ does not depend on the strategy. For the worked
strategy $\bar S=1.1904\,\tha$, so going from $r=0$ to $r=1$ costs $1000\times1.1904=1{,}190.4\,\mxnha$ and gains
$0.45\times1.1904=0.5357\,\tCha$:

\begin{table}[h]
\begin{tabular}{@{}cccc@{}}
\toprule
$r$ & $\bar V$ ($\mxnha$) & $\bar C$ ($\tCha$) & $\CVaR_{0.25}$ ($\mxnha$) \\
\midrule
$0$ & $11{,}278.0$ & $0.1087$ & $4{,}612.0$ \\
$0.5$ & $10{,}682.8$ & $0.3765$ & $4{,}113.1$ \\
$1$ & $10{,}087.6$ & $0.6444$ & $3{,}614.3$ \\
\bottomrule
\end{tabular}
\caption{Retention changes nothing in the simulation; it only moves stover between income and carbon.}
\end{table}

\subsection{Limiting cases and two caveats on grain and stover}

\begin{itemize}
\item $s=0$: $D=0$ and $Q=Q(0)$, so $L_t=0$ and no fertilizer is bought (test \code{tests/test_nsga_planner.py:85}).
\item $r=1$: no stover income; $r=0$: carbon comes from roots alone. Both objectives are monotone in $r$, and $\bar Y,\bar L$ are non-decreasing in $s$ (checked on the whole lattice in both worlds, no violation).
\item $\CVaR_{0.25}\le\bar V$ always (it is a mean of the lowest quarter); also checked on the lattice.
\item Grain is dry matter but SIAP prices are on a reported (about $14\%$ moisture) basis, so revenue is understated by about $14\%$ for maize; second order.
\item \finding{Default-world stover and roots are not capped like grain.} The grain-N rule (\code{wofost_rotations.py:312--313}) caps grain, and it binds in $88\%$
  of default maize seasons at $s=0$ and $69\%$ at $s=1$ (never in the calibrated world). Stover and roots are not capped, so default maize gives
  $15.96\,\tha$ of stover beside $6.34\,\tha$ of grain at $s=1$ (all maize seasons, 2006--2021). Stover income is $14{,}639\,\mxnha$,
  $92\%$ of the default pure-maize net of $15{,}945$ at $s=0$. Every default-world money and carbon figure is inflated by this. Documented, not changed.
\end{itemize}

\section{Dominance, the lattice, brute force and how fronts are scored}
\label{sec:dominance}

\subsection{Pareto dominance}

Write every strategy's objectives in minimization form $\mathbf f(x)=(f_1,\dots,f_M)$, where $f_i=\sigma_i\times$ natural value
(\code{minimized}, :226--228). Strategy $x$ \emph{dominates} $x'$ when it is no worse everywhere and strictly better somewhere:
\begin{equation}
\mathbf f(x)\prec\mathbf f(x')\iff \forall i:\ f_i(x)\le f_i(x')\ \ \wedge\ \ \exists i:\ f_i(x)<f_i(x').
\label{eq:dom}
\end{equation}
The \emph{Pareto front} is $\mathcal P=\{x\in\Omega:\ \nexists\,x'\in\Omega,\ \mathbf f(x')\prec\mathbf f(x)\}$. The code
(\code{pareto_front}, :231--244) drops rows with identical objective vectors, then keeps the first non-dominated level of
pymoo's \code{NonDominatedSorting}. Objectives are first rounded to $9$ decimals (\codeline{FRONT_DECIMALS}), for duplicates and for dominance; compared exactly,
floating-point noise of order $10^{-16}$ decided dominance (Section~\ref{sec:artifact}).

\subsection{The lattice and brute force}

A strategy lives on the lattice $\Omega=\{\omega\}\times\{0.025\,a\}_{a=0}^{50}\times\{0.05\,b\}_{b=0}^{20}$, so
\begin{equation}
|\Omega| = 4\times51\times21 = 4{,}284 .
\label{eq:lattice}
\end{equation}
\codeline{N_SCALE_STEPS = 50}, \codeline{RETENTION_STEPS = 20} (:64--65); \code{brute_force} (:258--263) scores every point:
$0.26$\,s in total on this machine, about $60\,\mu$s per strategy (timing varies by machine). The lattice makes the box bounds
ordinary values: real-coded NSGA reaches the bounds, where these nearly linear trade-offs have their optima, only slowly.
Brute force makes the exact front known, so the search can be graded.

\subsection{Why the five-objective front has 483 strategies (and the default one 1{,}377)}

\textbf{Dominance resistance.} If objective values were independent continuous draws, two random points are mutually
non-dominated with probability
\begin{equation}
P(\text{mutually non-dominated})=1-\frac{2}{2^{M}}:\quad M=2:\ 0.50,\quad M=3:\ 0.75,\quad M=5:\ 0.94 .
\label{eq:resist}
\end{equation}
The objectives here are far from independent, but the lesson holds: with five objectives, almost nothing dominates
anything. The calibrated front is $483/4{,}284=11.3\%$ of the lattice; the default front is $1{,}377/4{,}284=32.1\%$.
The $483$ split into three blocks (composition computed by running \code{pareto_front} on the brute-force table):

\begin{table}[h]
\begin{tabularx}{\linewidth}{@{}R{0.15\linewidth}R{0.14\linewidth}R{0.16\linewidth}Y@{}}
\toprule
Rotation & Count & $s$ values & Why it is non-dominated \\
\midrule
maize & $231=11\times21$ & $0,0.025,\dots,0.25$ (11 values, the first cell) & Each $s$ step buys $+0.00117$ ($\tha$) of yield at $-90.7$ ($\mxnha$) of income and $+3.27$ ($\kgNha$) of N loss, so no step dominates another; the 21 retentions trade carbon against income and bad-year income. \\
maize--barley & $231=11\times21$ & same & Dominated by pure maize in the four objectives without bad-year income (all $1{,}071$ of its strategies are). Survives only through bad-year income: $-1{,}343$ vs $-4{,}187$ ($\mxnha$) for maize at $s=0$, $r=0$. \\
maize--faba bean & $21=1\times21$ & $0$ & $s=0$ is genuinely best on income, N loss and bad-year income. (Before the float fix $s=0.025$ also survived here through rounding noise; Section~\ref{sec:artifact}.) \\
\bottomrule
\end{tabularx}
\caption{Composition of the calibrated five-objective exact front (483 strategies).}
\label{tab:front504}
\end{table}

Inside the first cell $[0,0.25]$ the calibrated yield rises only from $2.3008$ to $2.3125\,\tha$ ($+0.0117$, $+0.5\%$) for $32.7\,\kgNha$
more N lost and $907\,\mxnha$ less income. A $0.5\%$ slope is below the model's noise, so the size of the front is a metric artifact,
not decision-relevant structure.

\textbf{The concrete pair: maize versus maize--faba bean.}

\begin{table}[h]
\begin{tabular}{@{}lcc@{}}
\toprule
Objective (calibrated) & maize, $s=0.25$, $r=0$ & maize--faba bean, $s=0$, $r=0$ \\
\midrule
yield $\bar Y$ ($\tha$) & $\mathbf{2.3125}$ & $2.0420$ \\
net income $\bar V$ ($\mxnha$) & $4{,}988$ & $\mathbf{12{,}528}$ \\
N lost $\bar L$ ($\kgNha$) & $32.7$ & $\mathbf{0}$ \\
carbon input $\bar C$ ($\tCha$) & $\mathbf{0.1104}$ & $0.1087$ \\
bad-year income $\CVaR_{0.25}$ ($\mxnha$) & $-5{,}187$ & $\mathbf{5{,}862}$ \\
\bottomrule
\end{tabular}
\caption{Bold marks the winner on each objective. Each wins at least once, so by Equation~\eqref{eq:dom} neither dominates the other.}
\end{table}

Maize wins on yield (by $0.27\,\tha$) only as a tonnage gap: the faba bean partner yields $1.77\,\tha$ against maize's $2.32$, yet sells for about twice the
price per tonne. It also wins on carbon by a trivial $0.0017\,\tCha$. Remove yield and the maize and
maize--barley blocks collapse: of the $462$ non-faba members, $444$ become dominated, leaving $18$ (5 maize with $r\ge0.8$ and 13 maize--barley with $r\ge0.9$, kept by carbon input).
(\textbf{Correction}: the Markdown draft said $374$ dominated and $88$ remaining; I could not reproduce that.)
Dropping bad-year income instead gives the four-objective front of $252$ ($231$ maize and $21$ faba bean members, no maize--barley).

\textbf{How big is the front at a decision-relevant resolution?} (my computation; the resolutions are my illustrative choice, not repository settings.)
Round each objective to a grid $\varepsilon_i$ before testing dominance,
\begin{equation}
\begin{gathered}
\hat f_i=\varepsilon_i\,\operatorname{round}\!\big(f_i/\varepsilon_i\big),\\
\boldsymbol\varepsilon=\big(0.05\ \tha,\ 250\ \mxnha,\ 2\ \kgNha,\ 0.02\ \tCha,\ 250\ \mxnha\big),
\end{gathered}
\label{eq:eps}
\end{equation}
for (yield, income, N loss, carbon, bad-year income). The calibrated five-objective front then has $22$ strategies ($11$ maize, $9$ maize--faba bean, $2$ maize--barley); with the coarser $(0.1,\ 500,\ 5,\ 0.05,\ 500)$ it has $15$.
The $483$ is therefore mostly resolution below any difference one could act on.

\subsection{A floating-point artifact inside the front (found by review, fixed)}
\label{sec:artifact}

\finding{Before the fix, $21$ of the $504$ calibrated front members (and of the $273$-member four-objective front) were maize--faba bean at $s=0.025$,
one for each retention. They were non-dominated only because their yield exceeded that at $s=0$ by $4.4\times10^{-16}\,\tha$.}
The cause is Equation~\eqref{eq:interp}: for maize--faba bean, $X_0=X_1$ in exact arithmetic (yield and carbon do not respond),
but $(1-w)X_0+wX_1$ is evaluated in floating point (\code{nsga_planner.py:142}) and lands one unit in the last place above $X_0$.
Exact comparison in \code{pareto_front} then saw ``strictly better in yield'' and kept the point, although it is worse on income, N loss and bad-year income.

\textbf{The fix.} \code{pareto_front} (\code{nsga_planner.py:231--244}) now rounds the minimization-form objectives to $9$ decimals before dropping duplicates and before testing dominance,
and returns the rows with their unrounded values:
\begin{equation}
\hat f_i=10^{-9}\operatorname{round}\!\big(10^{9}f_i\big),\qquad x\prec x'\ \text{is tested on}\ \hat{\mathbf f}.
\label{eq:tol}
\end{equation}
(\codeline{FRONT_DECIMALS = 9}, :66; tests \code{test_rounding_error_does_not_keep_a_strategy_on_the_front} and \code{test_front_rows_keep_their_unrounded_values}.)
\textbf{What changed:} the calibrated five- and four-objective fronts are $483$ and $252$ instead of $504$ and $273$, the robust four-objective front is $1{,}727$ instead of $1{,}730$,
and the final-population counts of the two calibrated NSGA-III runs fall to $118$ and $77$ (noise ties removed). The three-objective front ($26$), the default-world fronts ($1{,}377$)
and the robust three-objective front ($738$) do not change.
\textbf{One pick changed:} the calibrated sustainability-first pick on the exact front was maize--faba bean at $s=0.025$ (maize N $=4\,\kgNha$), one of the noise survivors; it is now $s=0$
(N $=0$), with net income $11{,}457$ instead of $11{,}394\,\mxnha$ and N lost $0$ instead of about $2\,\kgNha$. The first review had said that no pick would change; that was wrong.

\subsection{IGD+ and the hypervolume ratio}

To grade a found set $\mathcal A$ against the exact front $\mathcal P$, \code{front_quality} (:329--345) first rescales both
by the reference front's range, so that $\mathcal P$ spans $[0,1]$ on every objective:
\begin{equation}
\tilde{\mathbf f}=\big(\mathbf f-\mathbf z^{\mathrm{id}}\big)\oslash\big(\mathbf z^{\mathrm{nad}}-\mathbf z^{\mathrm{id}}\big),\qquad
z^{\mathrm{id}}_i=\min_{\mathbf g\in\mathcal P}g_i,\quad z^{\mathrm{nad}}_i=\max_{\mathbf g\in\mathcal P}g_i ,
\label{eq:rescale}
\end{equation}
where $\oslash$ is element-wise division (the span is replaced by 1 if zero). Then, with found points $\mathbf f\in\mathcal A$, reference points $\mathbf g\in\mathcal P$ and
the hypervolume reference point $\mathbf z^{\mathrm{ref}}=1.1\cdot\mathbf 1$ (all rescaled),
\begin{align}
\mathrm{IGD}^+(\mathcal A,\mathcal P)&=\frac1{|\mathcal P|}\sum_{\mathbf g\in\mathcal P}\ \min_{\mathbf f\in\mathcal A}
 \sqrt{\sum_{i=1}^{M}\max\{0,\ f_i-g_i\}^2},
\label{eq:igd}\\
\mathrm{HVR}(\mathcal A)&=\frac{\mathrm{HV}(\mathcal A)}{\mathrm{HV}(\mathcal P)},\qquad
\mathrm{HV}(\mathcal X)=\operatorname{vol}_M\!\Big(\bigcup_{\mathbf f\in\mathcal X}\{\mathbf m:\ \mathbf f\le\mathbf m\le\mathbf z^{\mathrm{ref}}\}\Big).
\label{eq:hvr}
\end{align}
IGD+ (inverted generational distance plus; dimensionless after rescaling) is the mean, over the reference points, of how far the
nearest found point is from \emph{weakly dominating} that reference point: a found point at least as good on every objective
contributes $0$. The hypervolume $\mathrm{HV}$ is the $M$-dimensional volume $\operatorname{vol}_M$ of the region dominated by the set
and bounded by $\mathbf z^{\mathrm{ref}}$.

\textbf{Tiny hand example} ($M=2$: N loss in $\kgNha$ and net income in $\mxnha$). Reference front: $(0,1000)$, $(40,3000)$, $(80,5000)$.
Ideal and nadir in minimization form are $(0,-5000)$ and $(80,-1000)$, so the rescaled reference points are
$(0,1)$, $(0.5,0.5)$, $(1,0)$. Found set $\mathcal A$: the two end points $(0,1)$ and $(1,0)$.
\begin{equation}
\begin{aligned}
\mathrm{IGD}^+ &= \frac{0+\sqrt{0^2+\max\{0,1-0.5\}^2}+0}{3}=\frac{0.5}{3}=0.1667,\\
\mathrm{HV}(\mathcal P) &= \underbrace{0.5\times0.1}_{x\in[0,0.5]}+\underbrace{0.5\times0.6}_{x\in[0.5,1]}+\underbrace{0.1\times1.1}_{x\in[1,1.1]}=0.05+0.30+0.11=0.46,\\
\mathrm{HV}(\mathcal A) &= 1\times0.1+0.1\times1.1=0.21,\qquad \mathrm{HVR}=\frac{0.21}{0.46}=0.4565 .
\end{aligned}
\label{eq:toyigd}
\end{equation}
\code{front_quality} returns $0.16667$ and $0.45652$ on these data (run). If $\mathcal A$ is better than $\mathcal P$
the ratio exceeds $1$ and IGD+ is $0$ (e.g.\ adding $(20,3500)$ and $(40,4000)$ gives $\mathrm{HVR}=1.61$).

\textbf{What they reward and hide.} IGD+ rewards covering \emph{every} reference point and ignores surplus points; the hypervolume ratio is dominated by the largest region, so a thin missing tail can cost little.
Neither says whether a missing point matters for a decision, and both depend on the reference front, which here includes the $21$ noise points above.
\textbf{Naive baseline}: Table~\ref{tab:random} compares NSGA with a uniform random sample of the same number of lattice points.

\section{NSGA-II and NSGA-III in this code}
\label{sec:nsga}

\textbf{Plain language.} NSGA (Non-dominated Sorting Genetic Algorithm) keeps a population of candidate strategies. Each
generation it breeds children by recombining and perturbing parents, then keeps the best half of parents plus children,
where ``best'' means (a) lower in the dominance order and (b) spread out along the front. NSGA-II measures spread with a
neighbor-gap statistic (crowding distance); NSGA-III, used from four objectives, measures it by how many members sit near each of a fixed
set of reference directions. The code (\code{optimize}, :266--326) uses NSGA-II for $M\le3$ and NSGA-III for $M\ge4$ (:308, :315).

\subsection{Encoding}

A candidate is three genes: \code{rotation} (a categorical choice of 4), \code{n_step} $a\in\{0,\dots,50\}$ and \code{r_step}
$b\in\{0,\dots,20\}$ (integers), decoded by \code{decode} (:292--294) to
\begin{equation}
s=0.025\,a,\qquad r=0.05\,b,\qquad f_i=\sigma_i\,f_i^{\mathrm{nat}}(\omega,s,r)\quad(\text{minimized, :305}).
\label{eq:encode}
\end{equation}
Every evaluation calls \code{Planner.objectives} (about $60$--$120\,\mu$s), so a run of $150$ generations takes $3$--$15$\,s.

\subsection{Non-dominated sorting and ranks}

Level $\mathcal F_1$ is the set of non-dominated members; remove it and the non-dominated members of the rest form $\mathcal F_2$, and so on.
Equivalently, the rank of $x$ is the length of the longest dominance chain ending in $x$.

\begin{table}[h]
\begin{tabular}{@{}lcccccc@{}}
\toprule
Toy point & A & B & C & D & E & F \\
$(f_1,f_2)$ & $(1,5)$ & $(2,3)$ & $(4,1)$ & $(2,5)$ & $(3,4)$ & $(5,5)$ \\
\midrule
Dominated by & --- & --- & --- & A, B & B & A, B, C, D, E \\
Level & $\mathcal F_1$ & $\mathcal F_1$ & $\mathcal F_1$ & $\mathcal F_2$ & $\mathcal F_2$ & $\mathcal F_3$ \\
\bottomrule
\end{tabular}
\caption{A toy sort. pymoo's \code{NonDominatedSorting} returns the same levels (run).}
\end{table}

\subsection{Crowding distance}

Within one level, sort the members by each objective $i$. A member's crowding distance is the average, over objectives,
of the normalized gap between its two neighbors (boundary members get $+\infty$, so extremes are always kept):
\begin{equation}
\mathrm{cd}(n)=\frac1M\sum_{i=1}^{M}\frac{f_i^{(n^+)}-f_i^{(n^-)}}{f_i^{\max}-f_i^{\min}},
\label{eq:cd}
\end{equation}
where $n^-$ and $n^+$ are the neighbors of $n$ in the order sorted by $f_i$ and the $f_i^{\max/\min}$ are taken over the level.
\textbf{Toy}: points $(0,4)$, $(0.5,3)$, $(2,1)$, $(4,0)$, both ranges $=4$.
\begin{equation}
\begin{aligned}
\mathrm{cd}\big((0.5,3)\big)&=\tfrac12\Big(\tfrac{2-0}{4}+\tfrac{4-1}{4}\Big)=\tfrac12(0.5+0.75)=0.625,\\
\mathrm{cd}\big((2,1)\big)&=\tfrac12\Big(\tfrac{4-0.5}{4}+\tfrac{3-0}{4}\Big)=\tfrac12(0.875+0.75)=0.8125 .
\end{aligned}
\end{equation}
The less crowded $(2,1)$ is preferred; pymoo's \code{calc_crowding_distance} returns the same values (run).

\subsection{NSGA-II selection and survival (\texorpdfstring{$M\le3$}{M at most 3})}

\textbf{Binary tournament} (\code{binary_tournament} with \code{comp_by_dom_and_crowding}, :309--314): draw two individuals; if one
dominates the other it wins, otherwise the one with larger $\mathrm{cd}$ wins (random on ties). \textbf{Survival}
(\code{RankAndCrowding}): merge the $n_{\mathrm{pop}}$ parents and $n_{\mathrm{pop}}$ children ($2n_{\mathrm{pop}}$), fill the next generation level by level, and truncate
the level that does not fit by decreasing $\mathrm{cd}$. Population $100$.

\subsection{Variation on mixed variables}

\emph{Code} (pymoo $0.6.2$, \code{pymoo/core/mixed.py}, read in the installed source): the two integer genes are crossed together with
SBX and mutated with PM, both followed by rounding to the nearest integer; the rotation gene is crossed with uniform crossover and mutated by
random choice. Duplicates (exactly equal genes) are removed.

\begin{table}[h]
\begin{tabularx}{\linewidth}{@{}l l Y Y@{}}
\toprule
Gene & Type & Crossover & Mutation \\
\midrule
\code{n_step}, \code{r_step} & Integer & SBX, $\eta_c=15$, applied with probability $0.9$ per mating, then each gene crossed with probability $0.5$; rounded & PM, $\eta_m=20$, per individual $0.9$, per gene $0.5$; rounded \\
\code{rotation} & Choice (4) & uniform crossover: genes swapped with probability $0.5$ & redraw uniformly over the 4 options with probability $0.5$ (mutation probability $1.0$ per individual) \\
\bottomrule
\end{tabularx}
\caption{Operators and their defaults as found in the pymoo source. The NSGA-III defaults (\code{SBX(eta=30)}) are overridden by \code{MixedVariableMating}.}
\end{table}

\textbf{SBX} (simulated binary crossover, bounded form as implemented). For one gene with bounds $[y_{\min},y_{\max}]$, parents $y_1<y_2$,
gap $\Delta_y=y_2-y_1$ and a uniform draw $\zeta\in(0,1)$ shared by both children, the children $y_1',y_2'$ are
\begin{equation}
\begin{aligned}
\beta_L&=1+\tfrac{2(y_1-y_{\min})}{\Delta_y},\quad A_L=2-\beta_L^{-(\eta_c+1)},\quad
\beta_q=\begin{cases}(\zeta A_L)^{\frac1{\eta_c+1}}, & \zeta\le 1/A_L,\\[2pt] (2-\zeta A_L)^{-\frac1{\eta_c+1}}, & \text{otherwise},\end{cases}\\
y_1'&=\tfrac12\big[(y_1+y_2)-\beta_q\,\Delta_y\big],\qquad
y_2'=\tfrac12\big[(y_1+y_2)+\beta_q'\,\Delta_y\big],
\end{aligned}
\label{eq:sbx}
\end{equation}
where $\beta_q'$ is $\beta_q$ computed from $\beta_U=1+2(y_{\max}-y_2)/\Delta_y$ instead of $\beta_L$.
\textbf{Toy}: parents $y_1=10$, $y_2=30$ on $[0,50]$, $\zeta=0.3$, $\eta_c=15$: $\beta_L=2$, $A_L=1.99998$, $\beta_q=0.96858$, so $y_1'=10.314$ and
$y_2'=29.686$ (pymoo's own \code{cross_sbx} with that draw gives $10.3142$ and $29.6858$). After rounding the children are $10$ and $30$: the crossover
changed nothing. By simulation of pymoo's operator, with parents $20$ steps apart the children equal the parents after rounding with probability
$0.55$, and with parents $2$ steps apart with probability $0.999$. So once the population converges, SBX with $\eta_c=15$ plus rounding is nearly a no-op
and exploration comes from mutation.

\textbf{PM} (polynomial mutation). For a gene value $y$ with $\xi_1=(y-y_{\min})/(y_{\max}-y_{\min})$, $\xi_2=(y_{\max}-y)/(y_{\max}-y_{\min})$ and a uniform draw $\zeta$:
\begin{equation}
\xi_q=\begin{cases}\big[2\zeta+(1-2\zeta)(1-\xi_1)^{\eta_m+1}\big]^{\frac1{\eta_m+1}}-1, & \zeta\le\tfrac12,\\[3pt]
1-\big[2(1-\zeta)+2(\zeta-\tfrac12)(1-\xi_2)^{\eta_m+1}\big]^{\frac1{\eta_m+1}}, & \zeta>\tfrac12,\end{cases}
\qquad y'=y+\xi_q\,(y_{\max}-y_{\min}).
\label{eq:pm}
\end{equation}
\textbf{Toy}: $y=10$ on $[0,50]$, $\eta_m=20$, $\zeta=0.2$: $\xi_1=0.2$, $(0.8)^{21}=0.00922$, bracket $=0.4+0.6\times0.00922=0.40553$,
$\xi_q=0.40553^{1/21}-1=-0.04207$, $y'=10-0.04207\times50=7.897\to8$ (pymoo: $7.8966$). Over many draws from $y=10$ the rounded move has mean size
$2.2$ steps; $19\%$ of mutations round back to $10$ and $66\%$ stay within $\pm2$.

\subsection{NSGA-III: reference directions, association, niching}

\textbf{Das--Dennis reference directions.} With $M$ objectives and $p$ partitions per axis, the reference directions are the points of the unit simplex whose coordinates are multiples of $1/p$:
\begin{equation}
\mathcal W=\Big\{\boldsymbol\psi=\tfrac1p(k_1,\dots,k_M):\ k_i\in\mathbb N_0,\ \sum_{i=1}^M k_i=p\Big\},\qquad
H=|\mathcal W|=\binom{p+M-1}{M-1}.
\label{eq:dd}
\end{equation}
For example $M=3$, $p=3$ gives $\binom52=10$ directions, from $(0,0,1)$ through $(\tfrac13,\tfrac13,\tfrac13)$ to $(1,0,0)$. The three cases used here
(\code{optimize}, :316--317) are, with the population size defaulting to $H$:

\begin{table}[h]
\begin{tabular}{@{}lccccl@{}}
\toprule
Run & $M$ & $p$ & $H=\binom{p+M-1}{M-1}$ & Population & Where \\
\midrule
4 objectives (default) & 4 & 6 & $\binom93=84$ & 84 & \codeline{{4: 6, 5: 5}} at :316 \\
5 objectives (default) & 5 & 5 & $\binom94=126$ & 126 & same \\
robust, 4 objectives & 4 & 10 & $\binom{13}3=286$ & 286 & \codeline{--partitions 10} \\
\bottomrule
\end{tabular}
\caption{Reference directions (counts confirmed by calling pymoo's \code{get_reference_directions}).}
\end{table}

\textbf{Normalization, association, niching} (\code{ReferenceDirectionSurvival}, read in the installed pymoo). Each generation the
objectives are mapped to $\tilde f_i=(f_i-z_i^{\mathrm{id}})/(a_i-z_i^{\mathrm{id}})$, with $z^{\mathrm{id}}$ the running ideal point and $a_i$ the
intercepts of the hyperplane through the extreme points (falling back to the worst value on the first level). Each member is associated with the reference ray at
the smallest perpendicular distance
\begin{equation}
d_\perp(\tilde{\mathbf f},\boldsymbol\psi)=\Big\|\tilde{\mathbf f}-\frac{\boldsymbol\psi^{\!\top}\tilde{\mathbf f}}{\boldsymbol\psi^{\!\top}\boldsymbol\psi}\,\boldsymbol\psi\Big\|_2 ,\qquad
h(\tilde{\mathbf f})=\arg\min_{\boldsymbol\psi\in\mathcal W}d_\perp(\tilde{\mathbf f},\boldsymbol\psi),
\label{eq:perp}
\end{equation}
and each ray $\boldsymbol\psi_h$ has a niche count $\nu_h$ (members of the already accepted levels associated with it). The splitting level is filled by repeatedly
picking the niche with the smallest $\nu_h$ among those that still have candidates: if $\nu_h=0$ take the candidate with the smallest $d_\perp$, otherwise a random one.
\textbf{Toy} ($M=2$, $p=2$, so $H=3$ rays $(1,0)$, $(\tfrac12,\tfrac12)$, $(0,1)$; point $\tilde{\mathbf f}=(0.6,0.2)$, $|\tilde{\mathbf f}|^2=0.40$):
\begin{equation}
\begin{aligned}
d_\perp(\cdot,(1,0))&=\sqrt{0.40-0.6^2}=0.200,\qquad
d_\perp(\cdot,(\tfrac12,\tfrac12))=\sqrt{0.40-\big(\tfrac{0.8}{\sqrt2}\big)^2}=\sqrt{0.08}=0.283,\\
d_\perp(\cdot,(0,1))&=\sqrt{0.40-0.2^2}=0.600 ,
\end{aligned}
\end{equation}
so the point belongs to the first niche. Selection for mating is a random tournament (\code{comp_by_cv_then_random}: with no constraints it picks one of the
two at random), so all selection pressure comes from survival.

\textbf{Why NSGA-III from four objectives.} With many objectives most of a population is mutually non-dominated, so ranks stop discriminating and a
neighbor-gap statistic stops being meaningful. Evidence in this repository: in all five NSGA-III runs of Table~\ref{tab:bug} the final
population is almost entirely non-dominated ($84$, $126$ and $286$ members for the default and robust runs; $77$ and $118$ of $84$ and $126$ for the calibrated runs once noise ties are removed), and on the test fixture all $462$ strategies are mutually non-dominated
(\code{tests/test_nsga_planner.py:199--204}). That the replacement works better than NSGA-II would here is the standard argument (Deb and Jain, 2014); I did not
run NSGA-II with $M\ge4$. The normalization looked sound on the calibrated four-objective run (100 generations, seed 1): the ideal point matched the exact front's
exactly and the nadir was close ($36.6$ against $32.7\,\kgNha$ for N loss).

\subsection{The reporting bug that was fixed: \texorpdfstring{\code{result.X}}{result.X} versus the final population}

\emph{Code, before.} \code{optimize} returned pymoo's \code{result.X}. For NSGA-III, pymoo's \code{_set_optimum} sets this to
\code{survival.opt}: the first-level members that are, for some \emph{occupied} reference direction, the closest to that ray. Hence
\begin{equation}
\big|\texttt{result.X}\big|\ \le\ \Big|\big\{h(\tilde{\mathbf f}_n):\ n\in\text{merged population}\big\}\Big|
\label{eq:occ}
\end{equation}
(the number of occupied niches, i.e.\ reference rays that are the nearest ray of at least one member), regardless of how many non-dominated strategies the population holds. \emph{Code, now}
(:323--326): \code{result.pop} is decoded and reduced to its non-dominated set. Re-running the documented commands (seed $1$):

\begin{table}[h]
\begin{tabular}{@{}lrrrrrrrr@{}}
\toprule
 & & \multicolumn{2}{c}{\code{result.X} (before)} & \multicolumn{2}{c}{final population (now)} & \multicolumn{2}{c}{all evaluated} \\
\cmidrule(lr){3-4}\cmidrule(lr){5-6}\cmidrule(l){7-8}
Run (generations, $H$) & Exact & count & HVR & count & HVR & unique & HVR \\
\midrule
calibrated 4 obj (200, 84) & 252 & 18 & 0.937 & 77 & 0.981 & 1{,}485 & 1.000 \\
calibrated 5 obj (150, 126) & 483 & 16 & 0.888 & 118 & 0.985 & 1{,}624 & 1.000 \\
default 4 obj (200, 84) & 1{,}377 & 23 & 0.908 & 84 & 0.954 & 3{,}341 & 1.000 \\
default 5 obj (150, 126) & 1{,}377 & 22 & 0.925 & 126 & 0.965 & 3{,}248 & 0.999 \\
robust 4 obj (150, 286) & 1{,}727 & 45 & 0.927 & 286 & 0.977 & 4{,}246 & 1.000 \\
\bottomrule
\end{tabular}
\caption{Coverage by hypervolume ratio (Equation~\eqref{eq:hvr}) before and after the fix. ``All evaluated'': non-dominated set of every distinct lattice point the run evaluated.
The ``now'' columns equal the saved \code{run.json} values (e.g.\ $0.9850$, $0.9809$, $0.9544$).}
\label{tab:bug}
\end{table}

The count returned equals the number of occupied niches in every run ($18/84$, $16/126$, $23/84$, $22/126$, $45/286$; I read the niche of each final member),
so $4$--$8$ population members pile into each occupied niche. The search was never deficient: the final population already holds $95$--$98.5\%$ of the
hypervolume, and the archive of everything evaluated recovers essentially the whole front. What limits the returned front is the population cap
($77$--$286$ returned against fronts of $252$--$1{,}727$; the two calibrated counts are below the population size because noise ties are removed). IGD+ before $\to$ now: $0.016\to0.006$, $0.018\to0.005$, $0.048\to0.029$, $0.061\to0.030$, $0.037\to0.016$.


\textbf{NSGA-II on three objectives.} It is exact when the front fits in the population ($26$ of $26$, seeds 1--3, IGD+ $0$) and otherwise keeps
$0.965$--$0.973$ of the hypervolume (default, $1{,}377$ exact points, population $100$; IGD+ $0.023$--$0.025$).

\subsection{Is the search worth it at this size? A naive baseline}

\begin{table}[h]
\setlength{\tabcolsep}{5pt}
\begin{tabular}{@{}lrcccr@{}}
\toprule
Run & Evaluated & Returned & Archive & Random & Exact \\
\midrule
calibrated 4 obj & 1{,}485 (35\% of 4{,}284) & 0.981 & 1.000 & 0.9665 & 252 \\
calibrated 5 obj & 1{,}624 (38\%) & 0.985 & 1.000 & 0.9661 & 483 \\
default 4 obj & 3{,}341 (78\%) & 0.954 & 1.000 & 0.9983 & 1{,}377 \\
default 5 obj & 3{,}248 (76\%) & 0.965 & 0.999 & 0.9985 & 1{,}377 \\
robust 4 obj & 4{,}246 (99\%) & 0.977 & 1.000 & 0.9999 & 1{,}727 \\
\bottomrule
\end{tabular}
\caption{Hypervolume ratio for NSGA's \emph{returned} final population, NSGA's \emph{archive} (everything evaluated), and a uniform random sample of the same number of lattice points (``Random'': mean of 10 draws, my run); ``Exact'' is the exact front size.
The returned set is capped by the population size, a random sample's is not, so compare the archive with the random column.}
\label{tab:random}
\end{table}

NSGA beats random sampling at equal evaluation count only where the front is small (calibrated), and it visits $35$--$99\%$ of the lattice, so against brute force
(all $4{,}284$ in $0.26$\,s) it saves nothing. Its role here is a check of the machinery, which would earn its keep only in a much larger decision space.
\textbf{Why a lattice} (versus the earlier real-coded genes). I re-ran an emulation of the real-coded variant: on the test fixture (population and generations $40$) IGD+ was $0.098$, $0.172$, $0.133$ for seeds 1--3;
on the real data (population 100, 100 generations) $0.010$--$0.011$ (calibrated) and $0.018$--$0.020$ (default), against $0.000$ and $0.023$--$0.025$ on the lattice. The lattice matters at small budgets and is a wash on the real problem.

\section{The compromise picker \texorpdfstring{\code{choose}}{choose}: weighted Chebyshev}
\label{sec:choose}

\textbf{Plain language.} A front of hundreds of strategies is not a decision. \code{choose} (:348--358) returns the front member that is
closest to the ``utopia'' point (best value of every objective on that front) after each objective is rescaled to $[0,1]$ and weighted.
Distance is the \emph{largest} weighted shortfall, so the pick is the strategy whose worst weighted objective is least bad.

\subsection{The equation}

For a front with members $n=1,\dots,|\mathcal P|$ and minimization-form values $f_{ni}$,
\begin{equation}
\begin{gathered}
\tilde f_{ni}=\frac{f_{ni}-\min_{n'}f_{n'i}}{\max_{n'}f_{n'i}-\min_{n'}f_{n'i}}\in[0,1],\qquad
\lambda_i=\frac{v_i}{\sum_{i'}v_{i'}},\\
\gamma_n=\max_{i}\ \lambda_i\,\tilde f_{ni},\qquad
n^\ast=\arg\min_n\gamma_n .
\end{gathered}
\label{eq:cheb}
\end{equation}
Here $\tilde f_{ni}$ is a dimensionless shortfall ($0$ = best on the front, $1$ = worst), $v_i$ are the user's weights (:85--88), $\lambda_i$ their normalization and
$\gamma_n$ is the score (smaller is better). Code: the whole of Equation~\eqref{eq:cheb} is the single expression at :358, with the
span replaced by $1$ when an objective is constant on the front (:356). The ranges come from \emph{the front that is passed in}.

\textbf{Toy.} Front (yield, N loss) $=(3,100)$, $(2,50)$, $(1,0)$, weights $0.6$ (yield) and $0.4$ (N loss); yield is maximized, so $f_1=-3,-2,-1$:
\begin{equation}
\tilde{\mathbf f}=\begin{pmatrix}0&1\\0.5&0.5\\1&0\end{pmatrix},\quad
\gamma=\begin{pmatrix}\max(0.6\cdot0,\ 0.4\cdot1)\\ \max(0.6\cdot0.5,\ 0.4\cdot0.5)\\ \max(0.6\cdot1,\ 0.4\cdot0)\end{pmatrix}
=\begin{pmatrix}0.4\\0.3\\0.6\end{pmatrix}\ \Rightarrow\ n^\ast=2\ (\text{the middle one}).
\end{equation}
\code{choose} returns index $1$ (zero-based), as it should (also \code{tests/test_nsga_planner.py:290--295}).

\textbf{Assumptions.} A Chebyshev score can reach non-convex parts of a front, unlike a weighted sum, and the weights are importance judgments
($0.30$ carbon, $0.20$ bad-year income, $0.20$ N loss, $0.20$ income, $0.10$ yield for ``sustainability first''). But it uses only the single binding objective of each member.

\subsection{Why picks are unstable}

\begin{enumerate}
\item \textbf{Only the binding objective counts.} On the exact robust four-objective front ($1{,}727$ strategies), ``sustainability-first'' scores
  $0.1065$ at maize N $=128$, $0.1066$ at $124$, $0.1067$ at $120$ and $0.1082$ at $88\,\kgNha$, all bound by the carbon shortfall; $36$ members score within $0.01$ of the
  best, spanning $0$--$140\,\kgNha$. The dose is decided in the fourth decimal.
\item \textbf{Ties are broken by row order.} In ``profit-first'' on the same front, $547$ members tie at exactly the minimum score $0.2462$ because the same
  yield term binds for all of them, and \code{np.argmin} returns the first row (maize N $=4\,\kgNha$).
\item \textbf{Ranges come from the front itself.} The pick moves when the front is sampled instead of exact (Table~\ref{tab:picks}).
\item \textbf{Range normalization amplifies trivia.} On the calibrated exact front the yield range is $2.042$--$2.3125\,\tha$ and income $1{,}824$--$12{,}528\,\mxnha$:
\end{enumerate}

\begin{table}[h]
\setlength{\tabcolsep}{5pt}
\begin{tabular}{@{}llccccc@{}}
\toprule
Strategy & & income & yield & bad-year & N loss & carbon \\
 & weight $\lambda_i$ & $0.40$ & $0.30$ & $0.10$ & $0.10$ & $0.10$ \\
\midrule
maize, $s=0$, $r=0$ & $\tilde f_{ni}$ & $0.620$ & $0.043$ & $0.813$ & $0$ & $0.998$ \\
(income $5{,}895$) & $\lambda_i\tilde f_{ni}$ & $0.248$ & $0.013$ & $0.081$ & $0$ & $0.100$ \\
 & $\gamma_n=\max$ & \multicolumn{5}{l}{$\mathbf{0.248}$ (income binds)} \\
\midrule
maize--faba bean, $s=0$, $r=0$ & $\tilde f_{ni}$ & $0$ & $1.000$ & $0$ & $0$ & $1.000$ \\
(income $12{,}528$) & $\lambda_i\tilde f_{ni}$ & $0$ & $0.300$ & $0$ & $0$ & $0.100$ \\
 & $\gamma_n=\max$ & \multicolumn{5}{l}{$0.300$ (yield binds)} \\
\bottomrule
\end{tabular}
\caption{Why ``profit-first'' picks pure maize on the calibrated exact front. For maize--faba bean, $\tilde f_{n,\mathrm{yield}}=(2.3125-2.042)/(2.3125-2.042)=1$.}
\end{table}

``Profit-first'' gives up $1-5{,}895/12{,}528=53\%$ of attainable income to avoid a $0.26\,\tha$ tonnage gap that is mostly an artifact of averaging maize and faba bean tonnes.

\begin{table}[h]
\setlength{\tabcolsep}{4pt}
\begin{tabular}{@{}llll@{}}
\toprule
Front & Weights & Exact-front pick & NSGA-front pick \\
\midrule
calibrated (5 obj) & sustainability & mf, $N=0$, $r=0.90$ & mf, $N=0$, $r=0.95$ \\
 & profit & mz, $N=0$, $r=0$ & mz, $N=0$, $r=0$ \\
default (5 obj) & sustainability & mf, $N=136$, $r=1.00$ & mf, $N=124$, $r=1.00$ \\
 & profit & mf, $N=156$, $r=0.75$ & mf, $N=152$, $r=0.75$ \\
robust (4 obj) & sustainability & mf, $N=128$, $r=1.00$ & mf, $N=104$, $r=1.00$ \\
 & profit & mf, $N=4$, $r=0$ (first of 547 ties) & mf, $N=8$, $r=0.10$ \\
\bottomrule
\end{tabular}
\caption{Compromise picks; mf = maize--faba bean, mz = maize, $N$ = maize N in $\kgNha$. Exact fronts: the saved \code{brute_force_front.csv}; NSGA: the saved \code{front.csv}.}
\label{tab:picks}
\end{table}

\textbf{Takeaway.} Use the picks to see the trade-offs, not as findings; in a space this small choose from the brute-force front, and report the shortlist around the pick.

\section{The robust planner \texorpdfstring{\code{RobustPlanner}}{RobustPlanner}}
\label{sec:robust}

\textbf{Plain language.} The default and calibrated parameter sets disagree completely about what nitrogen buys, and the
default world's money is about $8\times$ larger. The robust planner scores a strategy in \emph{both} worlds, measures in each world how far it is from the best
attainable value of that objective as a share of that world's range, and keeps the \emph{worse} of the two. A strategy is robust if it is never far from the best in
either world (a minimax-regret flavour).

\subsection{The equations}

Let $f_i^{(\theta)}(x)$ be objective $i$ of strategy $x$ in world $\theta\in\{\mathrm{def},\mathrm{cal}\}$, in minimization form. Over the lattice $\Omega$ of that world,
\begin{equation}
\begin{aligned}
z_i^{(\theta)}&=\min_{x'\in\Omega}f_i^{(\theta)}(x'),
&\Delta_i^{(\theta)}&=\max_{x'\in\Omega}f_i^{(\theta)}(x')-z_i^{(\theta)},\\
\delta_i^{(\theta)}(x)&=\frac{f_i^{(\theta)}(x)-z_i^{(\theta)}}{\Delta_i^{(\theta)}},
&\rho_i(x)&=\max_{\theta}\ \delta_i^{(\theta)}(x).
\end{aligned}
\label{eq:robust}
\end{equation}
$\delta_i^{(\theta)}\in[0,1]$ is the shortfall (dimensionless; $0$ = the best attainable in that world, $1$ = the worst) and $\rho_i$ is the robust objective, always
minimized (named \code{shortfall_<objective>}). Code: $z$ is \code{self.best[world]} (:211) and $\Delta$ is \code{self.span[world]} (:212, set to $1$ if constant); the division is at :220--221 and the maximum over worlds at :222;
the lattice values come from \code{brute_force} per world (:210). The same search machinery then runs on $\rho_1,\dots,\rho_M$. Tests: \code{tests/test_nsga_planner.py:224--246}.

\subsection{A worked strategy in both worlds}

Strategy: maize--faba bean, $s=0.55$ (maize N $=88\,\kgNha$), $r=0$. World constants: net income best/span (def) $106{,}995.47$ / $105{,}689.30$ and (cal) $12{,}528.02$ / $13{,}454.44\,\mxnha$.

\begin{equation}
\begin{aligned}
\delta^{\mathrm{def}}_{\text{income}}&=\frac{106{,}995.47-95{,}902.38}{105{,}689.30}=\frac{11{,}093.09}{105{,}689.30}=0.1050,\\
\delta^{\mathrm{cal}}_{\text{income}}&=\frac{12{,}528.02-11{,}153.02}{13{,}454.44}=\frac{1{,}375.00}{13{,}454.44}=0.1022,\\
\delta^{\mathrm{def}}_{\text{N loss}}&=\frac{27.276-0}{98.940}=0.2757,\qquad
\delta^{\mathrm{cal}}_{\text{N loss}}=\frac{54.720-0}{192.178}=0.2847 .
\end{aligned}
\end{equation}

\begin{table}[h]
\begin{tabular}{@{}lrcrcc@{}}
\toprule
Objective & default value & $\delta^{\mathrm{def}}$ & calibrated value & $\delta^{\mathrm{cal}}$ & $\rho=\max$ (world) \\
\midrule
net income ($\mxnha$) & $95{,}902$ & $0.1050$ & $11{,}153$ & $0.1022$ & $0.1050$ (def) \\
N loss ($\kgNha$) & $27.3$ & $0.2757$ & $54.7$ & $0.2847$ & $0.2847$ (cal) \\
yield ($\tha$) & $7.12$ & $0.2330$ & $2.04$ & $0.7808$ & $0.7808$ (cal) \\
carbon input ($\tCha$) & $1.07$ & $0.9859$ & $0.109$ & $1.0000$ & $1.0000$ (cal) \\
bad-year income ($\mxnha$) & $87{,}750$ & $0.0090$ & $4{,}487$ & $0.0840$ & $0.0840$ (cal) \\
\bottomrule
\end{tabular}
\caption{The five shortfalls of the worked strategy (code output matches the hand values in the first two rows exactly).}
\end{table}

\subsection{Why range normalization, and its weaknesses}

The worlds' money differs about $8\times$ (best income $106{,}995$ versus $12{,}528\,\mxnha$), so raw regret in pesos would let the default world decide everything; dividing by each world's range
puts both on $[0,1]$. But:
\begin{itemize}
\item \textbf{Narrow ranges inflate trivia.} The calibrated yield span is $0.346\,\tha$ ($1.966$--$2.312$) against $6.78$ in the default world, so any change of rotation reads as a huge
  shortfall ($0.78$ wherever faba bean is in the rotation). Leave yield out of the robust objectives unless tonnage matters.
\item \textbf{Spans include dominated corners.} The calibrated N-loss span is $192\,\kgNha$ (maize at $s=1.25$, $200\,\kgNha$ dose), while the worst N loss on the exact front is $32.7$. N-loss shortfalls
  therefore depend on the arbitrary grid upper bound of $1.25$.
\item \textbf{The ideal point is per objective}, so it is not jointly attainable; and weighing the two worlds equally is a modelling choice, not a probability.
\end{itemize}

\subsection{The minimax dose, derived}

Fix the rotation (maize--faba bean) and $r=0$, and vary maize N $=N$. Let $\bar V^\ast_\theta$ be the best mean income in world $\theta$ and
$R_\theta(N)=\bar V^\ast_\theta-\bar V_\theta(N)$ the regret in pesos ($\mxnha$). Whatever the normalizer $\kappa_\theta$ (range, best income, or $1$), the income shortfall is $R_\theta/\kappa_\theta$.
\begin{itemize}
\item \textbf{Calibrated world:} yield does not respond, so income is exactly linear in the dose. The phase-average dose is $\mu N$ with $\mu=(1+\tfrac{40}{160})/2=0.625$, hence
  $R_{\mathrm{cal}}(N)=m_{\mathrm{cal}}N$ with $m_{\mathrm{cal}}=\pi_N\mu=25\times0.625=15.625\ \mxnkgN$ (check: at $88$, $15.625\times88=1{,}375.0$, as above).
\item \textbf{Default world:} income rises with the dose and is best at $N=200$, so the regret falls from $R_{\mathrm{def}}(0)$ to $0$ at an average slope $m_{\mathrm{def}}$ (cell slopes $85$ to $107$):
\end{itemize}
\begin{equation}
\begin{aligned}
R_{\mathrm{def}}(0)&=106{,}995.47-87{,}966.19=19{,}029.28\ \mxnha,\\
m_{\mathrm{def}}&=\frac{19{,}029.28}{200}=95.1\ \mxnkgN .
\end{aligned}
\end{equation}
The robust income shortfall is the larger of a falling and a rising line, so it is smallest where they cross:
\begin{equation}
\frac{R_{\mathrm{def}}(0)-m_{\mathrm{def}} N^\ast}{\kappa_{\mathrm{def}}}=\frac{m_{\mathrm{cal}} N^\ast}{\kappa_{\mathrm{cal}}}
\quad\Longrightarrow\quad
N^\ast=\frac{R_{\mathrm{def}}(0)}{m_{\mathrm{def}}+m_{\mathrm{cal}}\,\Lambda},\qquad \Lambda=\frac{\kappa_{\mathrm{def}}}{\kappa_{\mathrm{cal}}} .
\label{eq:minimax}
\end{equation}
Only the ratio $\Lambda$ matters: it is how many calibrated-world pesos one default-world peso is worth. With the repository's range normalization,
$\Lambda=105{,}689.30/13{,}454.44=7.855$:
\begin{equation}
N^\ast=\frac{19{,}029.28}{95.15+15.625\times7.855}=\frac{19{,}029.28}{95.15+122.73}=\frac{19{,}029.28}{217.88}=87.3\ \kgNha .
\end{equation}
Solving the same equation on the exact piecewise-linear curves gives $89.4$; on the lattice (4\,kg steps) the worst shortfall is $0.1084$ at $84$, $0.1050$ at $88$ and $0.1068$ at $92$, so the
minimax dose is $88$. \emph{Beyond $N^\ast$}, the calibrated income shortfall rises, and N loss rises with dose in both worlds, so those doses are dominated on (income, N loss): the two-objective robust front has $23$ strategies, all maize--faba bean
at $r=0$ with maize N from $0$ to $88\,\kgNha$ ($0,4,\dots,88$).

\subsection{The same dose for every rotation, and its dependence on how regret is defined}

\begin{table}[h]
\setlength{\tabcolsep}{5pt}
\begin{tabular}{@{}lcrrrr@{}}
\toprule
Rotation ($r=0$) & $N^\ast$ & cal.\ $\bar V$ & cal.\ $\bar L$ & def.\ $\bar V$ & def.\ $\bar L$ \\
\midrule
maize & $100$ & $3{,}488$ & $92.2$ & $36{,}375$ & $49.5$ \\
maize--barley & $84$ & $2{,}375$ & $54.4$ & $26{,}989$ & $28.6$ \\
maize--faba bean & $88$ & $11{,}153$ & $54.7$ & $95{,}902$ & $27.3$ \\
maize--barley--faba bean & $72$ & $7{,}905$ & $38.8$ & $70{,}339$ & $19.4$ \\
\bottomrule
\end{tabular}
\caption{Minimax maize N $N^\ast$ ($\kgNha$) on the lattice (range-normalized income shortfall) and the income $\bar V$ ($\mxnha$) and N lost $\bar L$ ($\kgNha$) there, in each world. Reproduced from the code; matches \code{NSGA_RESULTS.md} on every row. For maize--faba bean at $88$, income is $89\%$ (cal) and $90\%$ (def) of the best attainable.}
\end{table}

\begin{table}[h]
\setlength{\tabcolsep}{4pt}
\begin{tabular}{@{}lccccc@{}}
\toprule
Normalizer $\kappa_\theta$ & $\Lambda$ & closed form & exact crossing & lattice, mf & lattice, mz / mb / mbf \\
\midrule
range of income (the repo) & $7.855$ & $87.3$ & $89.4$ & $88$ & $100\ /\ 84\ /\ 72$ \\
share of the best income & $8.540$ & $83.2$ & $85.1$ & $84$ & $84\ /\ 60\ /\ 56$ \\
absolute pesos ($\kappa=1$) & $1$ & $171.8$ & $172.7$ & $176$ & $200\ /\ 200\ /\ 200$ \\
\bottomrule
\end{tabular}
\caption{Minimax maize N ($\kgNha$) by normalizer, through $\Lambda$ in Equation~\eqref{eq:minimax}; mf = maize--faba bean, mz = maize, mb = maize--barley, mbf = maize--barley--faba bean. In absolute pesos the other three rotations sit at the grid's upper bound ($200\,\kgNha$).}
\label{tab:regret}
\end{table}

\textbf{The hedge is a property of the normalizer.} Equation~\eqref{eq:minimax} makes it explicit: a larger $\Lambda$ (the calibrated world counts for more) lowers the dose, from $172$ at $\Lambda=1$ to $83$ at $\Lambda=8.5$.
Absolute pesos make the default world's larger money decide, so the ``hedge'' simply follows the default model. Prices matter as well: over N price $\in\{10,25,50\}$, faba bean price $\times\{0.5,0.75,1\}$ and stover price $\in\{0,500,1000\}\ \mxnt$,
the lattice minimax dose for maize--faba bean ranges from $48$ to $132\,\kgNha$ (maize alone: $100$ to $148$). The extremes are $48$ (N price $50$, faba bean at half price, stover $1{,}000$) and $132$ (N price $10$, faba bean at full price); 27 re-runs of the planner, my script.



The statement that survives is therefore: ``below the current $160$ and above $0$ under range normalization and these placeholder prices''. It is a hedge by construction, not a finding about nitrogen.

\section{Statistics and honesty check}
\label{sec:stats}

\subsection{The holdout (2022--2025)}

The planner is built twice from the same table: a design \code{Planner} on 2006--2021 and a holdout \code{Planner} on 2022--2025 (\code{load_planners}, :482--488).
The saved fronts carry the holdout values as \code{*_holdout} columns. For the two compromises picked from the calibrated NSGA front:

\begin{table}[h]
\begin{tabular}{@{}lcccccc@{}}
\toprule
 & & \multicolumn{2}{c}{design 2006--2021} & \multicolumn{2}{c}{holdout 2022--2025} \\
\cmidrule(lr){3-4}\cmidrule(l){5-6}
Pick (calibrated NSGA front) & $(s,r)$ & $\bar V$ & $\CVaR_{0.25}$ & $\bar V$ & $\CVaR$ ($q=1$) \\
\midrule
sustainability-first (maize--faba bean) & $(0,\,0.95)$ & $11{,}397$ & $4{,}914$ & $11{,}598$ & $8{,}843$ \\
profit-first (maize) & $(0,\,0)$ & $5{,}895$ & $-4{,}187$ & $6{,}158$ & $2{,}814$ \\
\bottomrule
\end{tabular}
\caption{Net income and bad-year income in $\mxnha$. With four holdout years $q=\lceil0.25\times4\rceil=1$, so the holdout ``bad-year'' value is the single worst year.}
\end{table}

The order of the two picks is kept. But the holdout does not test resilience: 2023 had $339.5$\,mm of season rain, below every design year (range $558$--$887$, mean $727$), yet the calibrated model still gives $1.99\,\tha$ of maize (design mean $2.30$), so
the calibrated model barely responds to drought. (The Markdown draft quoted $11{,}598$ against $6{,}086$: the sustainability-first pick is back at $11{,}598$ after the float fix, and the profit-first pick of the current front holds out at $6{,}158$.) The parameters were fit to 2006--2021 SIAP yields; 2022--2025 never selected the knobs but was looked at while the first calibration was designed (\code{CALIBRATION_RESULTS.md}).

\subsection{The CVaR over 16 years}

\begin{itemize}
\item \textbf{It is the mean of four numbers} (Equation~\eqref{eq:worked}): descriptive, with no confidence interval in the code. Illustration (my computation): an i.i.d.-year bootstrap of the worked strategy's 16 incomes ($20{,}000$ resamples, seed 0) gives a 5th--95th percentile range $[439,\ 8{,}220]\,\mxnha$ for $\CVaR_{0.25}$ (point value $4{,}113$) against $[8{,}598,\ 12{,}544]$ for the mean $\bar V$: the tail statistic is several times noisier.
\item \textbf{Bad years are frost years, not drought years, in the calibrated world.} Among the 16 pure-maize seasons at $s=0$, the point-biserial correlation of yield with a frost-ended season is $-0.70$ and with season rain $-0.02$
  (aggregated by year over all crops: $-0.65$ and $-0.19$). With $n=16$ the Fisher-$z$ 95\% interval for $-0.70$ is roughly $[-0.89,-0.31]$ (my arithmetic). \emph{Docs} (\code{CALIBRATION_RESULTS.md}): the POWER grid cell sits $\approx380$\,m above the fields and is
  about $2.5^\circ$C colder, which I did not re-derive. So ``bad-year income'' measures frost truncation and sowing date.
\end{itemize}

\subsection{Real data, placeholders, simulated, proxied}

\begin{table}[h]
\begin{tabularx}{\linewidth}{@{}R{0.22\linewidth}Y@{}}
\toprule
Category & Items \\
\midrule
Real data & POWER weather; IGBP-DIS and GLDAS soils; SIAP prices (median of three years) and the observed yields behind the calibration. But the faba bean price rests on $10$--$12$ ha of faba bean per year, and its latest value ($5{,}030$ in 2025) is $0.30\times$ the median used. \\
Placeholders & Costs $14{,}000$ / $9{,}000$ / $11{,}000\,\mxnha$, N price $25$, stover price $1{,}000$, carbon fraction $0.45$ (code comment: typical $0.40$--$0.47$). \\
Simulated (WOFOST) & Grain yield, N uptake, stover and root biomass, year by year, for chained seasons. \\
Proxied or accounting & \emph{Soil carbon}: an input flow, no decomposition, no stock. \emph{N loss}: accounting against hand-set fractions. \emph{Retention}: no feedback on yield, N or water. \emph{Prices}: constant. No N$_2$O or fertilizer-manufacturing emissions. \\
\bottomrule
\end{tabularx}
\end{table}

\subsection{Sensitivity of the exact-front picks}

Re-running \code{choose} on the exact fronts after one change at a time (``mf'' is maize--faba bean, ``mz'' maize, $N$ the maize dose in $\kgNha$; each cell reads sustainability-first / profit-first):

\begin{table}[h]
\begin{tabularx}{\linewidth}{@{}R{0.19\linewidth}YY@{}}
\toprule
Change & Calibrated (front size) & Default (front size) \\
\midrule
as shipped & mf, $N=0$, $r=0.90$ / mz, $N=0$, $r=0$ ($483$) & mf, $N=136$, $r=1.00$ / mf, $N=156$, $r=0.75$ ($1{,}377$) \\
N price $0$ & unchanged / mz, $N=40$, $r=0$ ($483$) & unchanged \\
N price $100$ & unchanged & unchanged / mf, $N=148$, $r=0.75$ \\
faba price $5{,}030$ (SIAP 2025) & mz, $N=0$, $r=0.55$ / mz, $N=0$, $r=0.35$ ($462$) & unchanged / mf, $N=200$, $r=0.35$ ($1{,}517$) \\
stover price $0$ & both mz, $N=0$, $r=1$ ($23$) & mz, $N=160$, $r=1$ / mf, $N=200$, $r=1$ ($102$) \\
carbon fraction $0.40$ & unchanged & unchanged (it cancels after range normalization) \\
\bottomrule
\end{tabularx}
\caption{One-at-a-time changes. The nitrogen response itself (calibrated versus default) changes every pick more than any price does; only a farm fertilizer trial can settle it.}
\label{tab:sens}
\end{table}

\section{Verification log, findings and what to keep in mind}
\label{sec:verify}
\small

\subsection{What I ran and what held}

\begin{itemize}
\item \codeline{python -m unittest tests.test_nsga_planner}: $35$ tests OK ($5.4$\,s).
\item Independent pandas recomputation from the raw \code{grid.csv} (10 strategies, both worlds, interpolated and node $s$): maximum relative difference $1.2\times10^{-16}$.
\item Worked strategy, node tables and leave-one-node-out error (Section~\ref{sec:objectives}, Tables~\ref{tab:years}, \ref{tab:nodes}, \ref{tab:loo}): hand and code agree.
\item Front sizes $483$, $1{,}377$, $252$, $26$, $738$, $1{,}727$ recomputed (after the float fix); the saved \code{brute_force_front.csv} sets equal the recomputed sets (calibrated, default).
\item NSGA re-runs, seed 1, documented commands: hypervolume ratios equal the saved \code{run.json} ($0.9850$, $0.9809$, $0.9544$, \ldots); grid hash \code{df8bfe11...} matches.
\item Minimax table, price sweep and three normalizers; picks on exact and NSGA fronts; sensitivity table; four-plan cross-table: reproduced.
\item Toys (sorting, crowding, SBX, PM, IGD+, HVR, perpendicular distance): checked against pymoo and \code{front_quality}.
\item Monotonicity over the whole lattice, both worlds: $\bar Y,\bar L$ non-decreasing in $s$; $\bar V$ non-increasing and $\bar C$ non-decreasing in $r$; $\CVaR\le\bar V$; no violation.
\item Random-sample baseline, bootstrap and correlations: run (Table~\ref{tab:random}, Section~\ref{sec:stats}).
\end{itemize}

\subsection{What I could not verify}

\begin{itemize}
\item \code{--build-grid}: needs PCSE and a NASA login; the saved grid was taken as given (its hash matches the runs; \code{grid.json} records commit \code{21e30ec}).
\item The SMAP correlation $r=0.87$ (needs a simulation) and which parameter set produced it.
\item The earlier real-coded implementation (never committed; I emulated it), and standard NSGA-II/III theory (quoted from memory; those papers are not in \code{Articles/}).
\item That the ``N lost'' figure is physically meaningful, and that default faba bean at $9.8\,\tha$ is plausible (it is not, against $1$--$2\,\tha$ typical in Mexico; \emph{docs}).
\end{itemize}

\subsection{Findings}

\begin{enumerate}
\item \finding{Floating-point survivors in the calibrated fronts}: $21$ strategies existed only through a $4.4\times10^{-16}\,\tha$ yield difference and one pick depended on them.
  \emph{Fixed} (Section~\ref{sec:artifact}).
\item \finding{Ties in \code{choose} are resolved by row order} (:358, \code{np.argmin}); $547$ robust-front members tie exactly in profit-first (Section~\ref{sec:choose}).
\item \finding{Default-world stover and roots are uncapped} (\code{wofost_rotations.py:312--318}): grain is capped by the grain-N rule but stover is not (Section~\ref{sec:objectives}). Open.
\item \emph{Minor documentation slips}, fixed in the repository documents: the results document said ``4--12 ha of faba bean'' (it is $10$--$12$ ha; $4$--$12$ is barley and faba bean together), and a code comment said ``16 of 84'' (now ``16--18 of 84'').
\end{enumerate}

\subsection{Things to keep in mind}

\begin{enumerate}
\item The calibrated ``zero N'' and the default ``85--100\% of the dose'' are properties of two parameter sets, not agronomy; the fertilizer decision is the one input the data cannot settle.
\item Do not read the $483$- and $1{,}377$-point fronts or the Chebyshev picks as findings: they come from tiny slopes ($0.5\%$ yield), dominance resistance and range normalization. At a decision-relevant resolution the front has about $15$--$22$ strategies.
\item The ``$50$--$130\,\kgNha$ hedge'' exists only under range-style regret and the placeholder prices (Table~\ref{tab:regret}).
\item Everything monetary in the default world is driven by uncapped stover and an implausible faba bean yield; the bad-year objective measures frost and sowing date, not drought.
\end{enumerate}

\textbf{A check you could run next.} Redo the resolution-rounded front of Section~\ref{sec:dominance} with the smallest income and yield differences a farmer would act on, and see which rotations and doses survive; then ask what interior grid nodes at $s=0.05,0.1,0.15$ (\code{--build-grid} with \code{GRID_SCALES} edited; needs PCSE) say about the response inside the first cell, where the calibrated front's maize strategies sit.

\textbf{Background references} (quoted from memory; not in the repository): Deb et al.\ (2002), \emph{IEEE TEVC} 6(2), NSGA-II; Deb and Jain (2014), \emph{IEEE TEVC} 18(4), NSGA-III; Das and Dennis (1998), \emph{SIAM J.\ Optim.} 8(3); Ishibuchi et al.\ (2015), IGD+; Rockafellar and Uryasev (2000), \emph{J.\ Risk} 2(3), CVaR.

\end{document}
